\chapter{妨害社会管理犯罪}

% Source PDF physical page 360; see audit/source_coverage.json.
\begin{EditionOpeningPanel}
重点内容：1.扰乱公共秩序罪；2.妨害司法罪；3.毒品犯罪。
\end{EditionOpeningPanel}

\section{扰乱公共秩序罪}

[问题]甲驾驶摩托车违章，警察正在处理其违章。甲抗拒处罚，不敢打警察，但将警察的摩托车砸坏。甲是否构成袭警罪？（砸警车案）

\subsection*{一、妨害公务类犯罪}

\subsubsection*{（一）妨害公务罪}

第277条第1款以暴力、威胁方法阻碍国家机关工作人员依法执行职务的，处三年以下有期徒刑、拘役、管制或者罚金。

\InsertDiagram{figures/manual_table_p360_1.tex}

% Source PDF physical page 361; see audit/source_coverage.json.
\subsubsection*{（二）袭警罪}

第277条第5款暴力袭击正在依法执行职务的人民警察的，处三年以下有期徒刑、拘役或者管制；使用枪支、管制刀具，或者以驾驶机动车撞击等手段，严重危及其人身安全的，处三年以上七年以下有期徒刑。［《刑法修正案（十一）》增设]\footnote{2025年1月18日最高人民法院、最高人民检察院《关于办理袭警刑事案件适用法律若干问题的解释》。}

1.行为对象：正在执行公务的警察。

（1）这里的“警察”，包括公安机关、国家安全机关、监狱等部门的警察和法院、检察院的司法警察。

（2）这里的“警察”，不包括正在依法配合警察执行职务的警务辅助人员。对其暴力袭击，有可能构成妨害公务罪。同时暴力袭击正在依法执行职务的警察和配合警察执行职务的警务辅助人员，按照袭警罪论处。（2025年试题）

（3）正在执行公务。例如，甲驾驶摩托车违章，警察处理完毕后，甲认为警察处理过于严厉，端了警察一脚后驾车逃离。甲不构成袭警罪。

[注意］警察在非工作时间遇有紧急情况依法履行职责，暴力袭击该警察，构成袭警罪。

（4）警察执法活动存在严重过错的，对行为人一般不作为犯罪处理；执法活动存在较大过错，袭击行为暴力程度较轻、危害不大的，可以不作为犯罪处理。

2.行为方式：暴力袭击。

（1）暴力袭击的方式：①实施撕咬、掌捆、踢打、抱摔、投掷物品等行为，造成轻微伤以上后果的；②实施打砸、毁坏、抢夺警察乘坐的车辆、使用的警械等行为，\CJKunderwave{足以危及}\CJKunderwave{人身安全}的。

（2）与警察发生轻微肢体冲突，或者为摆脱抓捕、约束实施甩手、挣脱、蹬腿等一般性抗拒行为，危害不大的，或者仅实施辱骂、讽刺等言语攻击行为的，不属于这里的“暴力袭击”。

（3）行为方式不包括以暴力相威胁。以暴力相威胁，可以构成妨害公务罪。

\subsection*{二、关于公文、证件、印章犯罪}

\subsubsection*{(一）伪造、变造、买卖国家机关公文、证件、印章罪（第 280 条)}

1.保护法益：国家机关公文、证件、印章的公共信用。

（1）实施本罪，不要求存在真实对应的国家机关及其公文、证件、印章，只要求足以使一般人相信有这个国家机关及其公文、证件、印章，因为这种行为侵害了国家机关公文、证件、印章的公共信用。例如，狗蛋伪造“卫生部”的公文。虽然我国目前没有“卫生部”，有“卫生健康委员会”（简称卫健委），但是，有些老百姓会误以为有“卫生部”。狗蛋的行为侵害了卫健委的公文的公共信用，属于伪造国家机关公文。

（2）成立本罪，要求有非法使用目的。如此，才可能侵害国家机关公文、证件、印章的公共信用。

例1，狗蛋出于好玩，伪造一个国家机关印章，然后放抽屉里，没有使用目的，不构成本罪。

% Source PDF physical page 362; see audit/source_coverage.json.
例2，林业部门应当为狗蛋发放林权证书，故意不发放。无奈，狗蛋伪造一个林权证书，抵押给小芳，向小芳借款。狗蛋没有非法使用目的，不构成本罪。

例3，狗蛋屡次相亲失败，女方嫌其不够成熟。为此，狗蛋伪造户口本，将自已的年龄改大，出示给女方看，以显示自己成熟。狗蛋的这种使用方式，不属于户口本的正常用途，不构成本罪。

2.伪造

（1）伪造包括有形伪造与无形伪造。有形伪造，是指没有制作权限的人，冒用国家机关名义制作公文、证件。无形伪造，是指有制作权限的人，擅自以国家机关的名义制作与事实不符的公文、证件；如果未经授权，擅自制作与事实相符的公文、证件，不属于伪造。

注意：本罪可以由间接正犯构成。例如，狗蛋提供虚假材料，欺骗国家机关工作人员，使其开具了与事实不符的证件。狗蛋构成伪造国家机关证件罪的间接正犯。（2023年试题）

（2）伪造印章，包括伪造印体，如私刻公章，也包括伪造印影。

例1，狗蛋用红笔描绘公章图案，属于伪造印章。不过，仅干这件事，尚不构成犯罪。狗蛋在某个文件上，描绘公章图案，既构成伪造公文，也构成伪造印章，构成伪造国家机关公文、印章罪。

例2，小芳负责保管某国家机关印章。狗蛋找到小芳，拿着“狗蛋系某法院执行局法官”的虚假证明，请小芳盖章。小芳明知内容有假而照办。虽然印章真实，但内容虚假，属于伪造国家机关公文。

3.变造

（1）变造是指对真实的国家机关公文、证件、印章进行修改，改变其非本质内容的行为。如果改变了公文、证件、印章的本质内容，则应认定为伪造。

（2）变造只能针对现实中的原件进行修改。伪造可以针对原件的电子图片进行修改。

例如，狗蛋将一份真实的公文扫描到电脑里，生成电子图片，然后修改，再用彩色打印机打印出来。这个过程属于伪造，而非变造。此时，无论修改的程度大小，都属于伪造。同理，狗蛋将一份真实公文的复印件扫描到电脑里，继续上述操作，也属于伪造。也即，伪造可以针对复印件进行修改。

4.买卖

（1）买卖的对象，包括真实的公文、证件、印章，也包括伪造、变造的公文、证件、印章。这是因为，买卖伪造的公文、证件、印章，也会侵害国家机关公文、证件、印章的公共信用和证明信用。\CJKunderwave{注意}：不包括民事判决书。

（2）买卖的方式。卖，包括先买进后卖出，也包括单纯的卖出。\CJKunderwave{买，包括为卖出而买}\CJKunderwave{进}，\CJKunderwave{也包括为自己使用而买进}。

[注意］伪造、变造、买卖机动车行驶证、登记证书，机动车人户、过户、验证的证明文件，构成本罪。但是，伪造、变造、买卖民用机动车号牌，不能认定为本罪。\footnote{参见2009年最高人民法院研究室《关于伪造、变造、买卖民用机动车号牌行为能否以伪造、变造、买卖国家机关证件罪定罪处罚问题的请示》的答复。其中理由：第280条的伪造、变造、买卖国家机关证件罪的法定刑高于第281条的非法生产、买卖警用装备罪的法定刑。若伪造民用机动车号牌定前罪，而伪造警用机动车牌定后罪，就导致伪造民用机动车号牌的处罚重于伪造警用机动车号牌。这显然不合理。}如果反复伪造、买卖民用机动车号牌，可定非法经营罪。

% Source PDF physical page 363; see audit/source_coverage.json.
5.罪数

（1）实施本罪后，又利用该公文、证件、印章实施诈骗罪，具有牵连关系的，择一重罪论处。

（2）非法经营罪有种行为类型：“买卖进出口许可证、进出口原产地证明以及其他法律、行政法规规定的经营许可证或者批准文件的行为”。该行为类型与买卖国家机关公文、证件罪形成想象竞合犯，择一重罪论处。

\subsubsection*{（二）伪造公司、企业、事业单位、人民团体印章罪（第 280 条)}

1.这里的公司、企业、事业单位、人民团体没有所有制的限制。例如，既包括公立医院，也包括私立医院。

2.行为对象只包括印章，不包括公文、证件，也即伪造公司、企业、事业单位、人民团体的公文、证件，不构成犯罪。

3.行为方式只包括伪造，不包括变造、买卖，也即变造、买卖公司、企业、事业单位、人民团体印章，不构成犯罪。变造、买卖这些单位的公文、证件，也不构成犯罪。

4.司法解释：“对于伪造高等院校印章制作学历、学位证明的行为，应当以伪造事业单位印章罪论处。明知是伪造高等院校印章制作的学历、学位证明而贩卖的，以伪造事业单位印章罪的共犯论处。”

\subsubsection*{(三）伪造、变造、买卖身份证件罪(第 280 条)}

1.身份证件，不仅包括居民身份证，还包括护照、社会保障卡、驾驶证等依法可以用于证明身份的证件。

[注意］国家机关制作的，仅供内部使用，用以证明身份的证件，例如工作证、出入证，不是这里的身份证件。

2.行为方式，包括伪造、变造、买卖。买卖：第一，买卖的对象，包括真实的身份证件，也包括伪造、变造的身份证件。第二，买卖的方式。卖，包括先买进后卖出，也包括单纯的卖出。\CJKunderwave{买，包括为卖出而买进，也包括为自己使用而买进}。

例1，清洁工甲拾得乙的身份证，卖给丙，构成本罪。

例2，甲将自已的身份信息告诉乙，让乙伪造，并支付办证费用。甲构成买卖身份证件罪。乙构成伪造身份证件罪。

\subsubsection*{（四）盗窃、抢夺、毁灭国家机关公文、证件、印章罪（第 280 条)}

1.行为对象是真实的公文、证件、印章。注意：买卖国家机关公文、证件、印章罪的对象可以是虚假的公文、证件、印章。

2.刑法规定了“盗窃、抢夺武装部队公文、证件、印章罪”（第375条）。若毁灭武装部队的公文、证件、印章的，可以定毁灭国家机关公文、证件、印章罪。（2014年第51题）

3.抢劫国家机关公文、证件、印章，可以定抢夺国家机关公文、证件、印章罪，因为抢劫可以包容评价为抢夺。

\subsubsection*{（五）冒名顶替罪}

1.行为方式：盗用、冒用他人身份+顶替三种事项。（上大学的资格，公务员录用资格，就业安置待遇）

% Source PDF physical page 364; see audit/source_coverage.json.
2.被顶替者即使同意，顶替者也构成本罪。这是因为，国之公器，不允许私相授受。

3.组织、指使他人实施冒名顶替，也构成冒名顶替罪，以正犯论处。

\subsection*{三、“小混混”型犯罪}

\subsubsection*{（一）招摇撞骗罪（第279条）}

这是指，冒充国家机关工作人员招摇撞骗的行为。本罪的保护法益是国家机关工作人员的公众信赖感（即官员的公众形象）。因此，冒充国家机关工作人员学雷锋、做好事，不构成犯罪。

1.行为方式。

（1）冒充的国家机关工作人员：不包括国有企业、事业单位人员，例如国企领导、大学校长；不包括高干子弟、烈士子女、战斗英雄和劳模。冒充已被撤销建制的国家机关的某个工作人员，足以使对方信以为真的，也可能成立本罪。

（2）冒充的情形：一是非国家机关工作人员冒充国家机关工作人员（包括离职的国家机关工作人员冒充在职的国家机关工作人员）；二是此种国家机关工作人员冒充彼种国家机关工作人员，如行政机关工作人员冒充司法机关工作人员；三是职务低的国家机关工作人员冒充职务高的国家机关工作人员，当然也包括相反情形（2019年试题）。

（3）招摇撞骗，要求有“欺骗”行为。这是指，采取欺骗手段，使对方产生认识错误，基于此骗取对方付出一些东西，包括感情、服务、财物等。

例1，甲冒充警察抢劫，不构成招摇撞骗罪，而构成抢劫罪，适用法定刑升格条件“冒充军警人员抢劫”。

例2，甲穿上林业部门人员的制服，冒充林业部门人员，盗伐路边林木，路人以为甲是林业部门人员，没有阻止。甲不构成招摇撞骗罪，因为甲虽然欺骗路人，但路人并没有付出什么东西，也即甲没有骗取路人的财物。（2023年试题）

[注意］成立本罪，并不要求骗钱，可以欺骗职务、资格、服务、感情等。当骗钱的时候，同时构成诈骗罪，想象竞合，择一重罪论处。诈骗罪与招摇撞骗罪不是法条竞合关系，因为触犯招摇撞骗罪并不必然触犯诈骗罪，例如，冒充官员招摇撞骗，骗取职务，只构成招摇撞骗罪，不构成诈骗罪。（法考界的多数说，2025年试题）

2.罪数问题。（2013年第10题）

（1）冒充军警人员抢劫的，属于抢劫罪的法定加重情节。

（2）冒充警察招摇撞骗，定招摇撞骗罪，并从重处罚。

（3）冒充军人（包括武警）招摇撞骗，定冒充军人招摇撞骗罪。

（4）冒充国家执法人员执法，以此骗取钱财，触犯了招摇撞骗罪和诈骗罪，按想象竞合犯处理。冒充国家执法人员执法，既欺骗又恐吓被害人，让被害人交付钱财，就触犯了招摇撞骗罪、诈骗罪、敲诈勒索罪，按照想象竞合犯处理。例如，冒充警察抓赌抓，就是这种情形。

\subsubsection*{（二）聚众斗殴罪（第292条）}

1.聚众斗殴是指以聚众的形式相互斗殴。（2013年第62题）

（1）相互斗殴，是指双方都出于侵害对方的意图而相互攻击，同时默示承诺对方对自已的侵害。如果一方有侵害意图，另一方没有侵害意图和行为，则实施侵害的一方属于故

% Source PDF physical page 365; see audit/source_coverage.json.
意伤害，可构成故意伤害罪，另一方可以正当防卫。

[提示]基于被害人承诺原理，聚众相互斗殴中，故意伤害致人轻伤，不构成故意伤害罪，故意伤害致人重伤、故意杀人，则构成故意伤害罪、故意杀人罪。

（2）聚众，是指要求多人参与，但不要求双方都必须3人以上。

例1，甲与乙均有相互侵害的意图，甲带领3人前往参加，乙一人参加，属于聚众斗殴。

例2，甲有侵害乙的意图，乙没有侵害甲的意图，甲带领3人前往，乙一人迎战。甲方属于故意伤害，乙属于正当防卫。

例3，甲有侵害乙的意图，乙没有侵害甲的意图，甲一人前往，乙带领3人迎战。甲属于故意伤害，乙方属于正当防卫。

［提示］聚众的“众”，不要求达到责任年龄、具有责任能力。

2.本罪只处罚首要分子和积极参加者。

3.第292条第2款规定：聚众斗殴过失致人重伤、死亡的，以故意伤害罪、故意杀人罪定罪处罚。（2014年第20题）

（1）该规定是法律拟制，是指过失致人重伤、死亡，也定故意伤害罪、故意杀人罪。若故意重伤或杀人，理所当然也定故意伤害罪、故意杀人罪。

（2）该规定是针对两种人而言的：一是直接造成重伤、死亡的斗殴者；二是首要分子。其他参与者对此不负责任。在不能查明死亡是由谁造成时，只让首要分子负责。（2022年试题）

（3）致人重伤、死亡，既包括一方导致另一方人员重伤、死亡，也包括一方导致本方人员重伤、死亡。

\subsubsection*{（三）寻衅滋事罪}

第293条有下列寻滋事行为之一，破坏社会秩序的，处五年以下有期徒刑、拘役或者管制：

\subsubsection*{（一）随意殴打他人，情节恶劣的；}

\subsubsection*{（二）追逐、拦截、辱骂、恐吓他人，情节恶劣的；}

\subsubsection*{（三）强拿硬要或者任意损毁、占用公私财物，情节严重的；}

\subsubsection*{（四）在公共场所起哄闹事，造成公共场所秩序严重混乱的。}

纠集他人多次实施前款行为，严重破坏社会秩序的，处五年以上十年以下有期徒刑，可以并处罚金。

1.行为类型。关于第四种行为类型，有关司法解释规定：\footnote{2013年9月10日最高人民法院、最高人民检察院《关于办理利用信息网络实施诽谤等刑事案件适用法律若干问题的解释》第5条第2款。}编造虚假信息，或者明知是编造的虚假信息，在信息网络上散布，或者组织、指使人员在信息网络上散布，起哄闹事，造成公共秩序严重混乱的，依照第293条第1款第4项的规定，以寻滋事罪论处。

2.罪数问题。旧观点认为，本罪与故意伤害罪、抢劫罪、敲诈勒索罪的区分，是对立排斥的区分。新观点及司法解释认为，二者可以想象竞合，然后择一重罪论处。\footnote{参见张明楷：《寻滋事罪探究（上、下篇）》，载《政治与法律》2008年第1期、第2期。}例如，“随意殴打他人，情节恶劣的”，与故意伤害罪可以想象竞合。（2015年第58题）

% Source PDF physical page 366; see audit/source_coverage.json.
\subsubsection*{（四）聚众扰乱社会秩序罪（第290条）}

1.聚众：是指组织聚集众人。如果未经组织，民众自发围观汇集，不是聚众。

2.行为方式是聚众，既可采用暴力手段，也可采用非暴力手段。如果是暴力手段，触犯其他犯罪，数罪并罚。例如，聚众时使用暴力故意致人重伤、死亡，还需定故意伤害罪、故意杀人罪。

[提示］假装跳楼、跳桥，吸引观众注意等“跳楼秀”，不构成本罪。

\subsubsection*{（五）高空抛物罪}

第291条之二从建筑物或者其他高空抛掷物品，情节严重的，处一年以下有期徒刑、拘役或者管制，并处或者单处罚金。

有前款行为，同时构成其他犯罪的，依照处罚较重的规定定罪处罚。［《刑法修正案

\subsubsection*{（十一）》增设]}

1.行为方式。

（1）从山上、树上、巨轮上、热气球上抛物，属于高空抛物。

（2）明知井下有人作业，仍从井口向下扔东西，属于高空抛物。也即“高空”是个相对概念。

（3）从地面向上抛物，使物品从高空坠落，不属于高空抛物。也即高空抛物，要求从“高空”向下抛物。

（4）多高算“高”空？高度是二层楼或3米左右。

2.情节犯。成立本罪，要求“情节严重”。对此需要综合考虑物品重量、数量、高度、危险程度等。成立本罪，不要求实害结果。

3.主观要件是故意。过失导致物品坠落，不构成本罪。

4.罪数问题。例1，甲向楼下有人行走的地方，扔个酒瓶子。这种行为没有危及多数人的可能性，没有危害公共安全，不构成以危险方法危害公共安全罪，而构成高空抛物罪。如果不慎砸死人，同时构成过失致人死亡罪，想象竞合，择一重罪论处。如果故意想砸死人，也砸死了人，同时构成故意杀人罪，想象竞合，择一重罪论处。\footnote{2019年的司法解释认为，“高空扔酒瓶子”，构成以危险方法危害公共安全罪。这是未增设“高空抛物罪”时的权宜之计。2021年《刑法修正案（十一）》增设“高空抛物罪”后，该司法解释规定便失效了。参见张明楷：《刑法学》（第六版），法律出版社2021年版，第108页。}（2026年试题）

例2，甲向楼下有人行走的地方，抛洒硫酸。这种行为危害了公共安全，同时构成以危险方法危害公共安全罪和高空抛物罪，想象竞合，择一重罪论处。

例3，甲向楼下有人行走的地方，扔燃烧的蜂窝煤球，有引发火灾危险。这种行为危害了公共安全，同时构成放火罪和高空抛物罪，想象竞合，择一重罪论处。这种行为也构成以危险方法危害公共安全罪，但在构成放火罪的情况下，优先定放火罪。（2022年试题）

\subsection*{四、“大哥”型犯罪}

\subsubsection*{（一）组织、领导、参加黑社会性质组织罪（第294条)}

1.“黑社会性质组织”。（1）认定标准：第一，具有稳定组织；第二，具有经济实力；第三，手段具有非法性、破坏性；第四，形成非法控制。\CJKunderwave{注意}：不要求必须有“政治保护伞”。（2019年试题）

% Source PDF physical page 367; see audit/source_coverage.json.
（2）主观认识。例如，甲明知参加了一个从事非法活动的组织，但不认为其是黑社会性质组织。法官认定该组织是黑社会性质组织。甲具有“参加黑社会性质组织罪”的犯罪故意。也即，甲只要认识到该组织是从事非法活动的组织即可，至于其是不是黑社会性质组织，则是法官认定的事情。（2024年试题）

2.责任界定。对于黑社会性质组织的组织者、领导者，应当按照其所组织、领导的\CJKunderwave{黑社会性质组织所犯的全部罪行}处罚。所谓黑社会性质组织所犯的全部罪行，不是指黑社会性质组织成员所犯的全部罪行，而是组织、领导者所组织、发动、指挥的全部罪行。成员个人所犯罪行，不属于黑社会性质组织所犯罪行。（2018年试题）

3.罪数。本罪的行为是组织、领导、参加黑社会性质组织。如果组织、领导、参加后，又实施了其他犯罪，数罪并罚。例如，组织黑社会性质组织后，又利用该组织实施赌博罪、组织卖淫罪、贩卖毒品罪等，应数罪并罚。

\subsubsection*{（二）赌博犯罪（第303条）}

1.赌博罪：以营利为目的，聚众赌博或者以赌博为业的行为。赌博，是指以偶然因素定输赢。

2.开设赌场罪：以营利为目的，设立专门用于赌博的场所，提供赌具的行为。常见方式：通过赌博网站、赌博机等方式开设赌场。

（1）如何区分网络游戏与网络赌博？标准：单向消费还是反向变现。

第一，服务商只提供单向消费，属于组织网络游戏。

第二，服务商提供反向变现，属于组织网络赌博，构成开设赌场罪。反向变现，是指玩家可以将赢取的积分、筹码、奖品等兑换为现金、货币。服务商提供这种变现渠道，或者默许玩家通过第三方进行变现，构成开设赌场罪。（2025年试题）

（2）最高人民法院、最高人民检察院指导性案例的要旨。

第一，以营利为目的，通过邀请人员加人微信群，利用微信群进行控制管理，以抢红包方式进行赌博，在一段时间内持续组织赌博活动的，以开设赌场罪论处。（最高人民法院第106号指导案例）（2021年试题）

第二，以营利为目的，通过邀请人员加人微信群的方式招揽赌客，根据竞猜游戏网站的开奖结果等方式进行赌博，设定赌博规则，利用微信群进行控制管理，在一段时间内持续组织网络赌博活动的，以开设赌场罪论处。（最高人民法院第105号指导案例）

（3）境外赌博网站的犯罪地

包括：网站服务器所在地，网络接人地，网站建立者、管理者所在地，参赌者所在地。例如，缅甸人甲在缅甸建立赌博网站，中国境内的人在该网站参赌，那么中国境内属于该开设赌场罪的犯罪地。中国公民丙为该网站提供帮助，构成开设赌场罪的帮助犯。（2021年试题）

3.组织参与国（境）外赌博罪：组织中国公民参与国（境）外赌博，数额巨大或者有其他严重情节的行为。［《刑法修正案（十一）》增设]

本罪包括情形：（1）组织中国公民赴国（境）外参与赌博；（2）组织中国公民在境内参与境外网络赌博；（3）组织身居境外的中国公民在境外参与赌博。

\subsubsection*{（三）虚假物质、信息犯罪}

1.投放虚假危险物质罪（第291条之一）

这是指，投放虚假的爆炸性、毒害性、放射性、传染病病原体等物质，严重扰乱社会秩序的行为。

% Source PDF physical page 368; see audit/source_coverage.json.
本罪只是扰乱社会秩序，没有危害公共安全。如果投放真实的危险物质，危害公共安全，则构成投放危险物质罪（第114条、第115条）。

2.编造、故意传播虚假恐怖信息罪（第291条之一）

这是指，编造爆炸威胁、生化威胁、放射威胁等恐怖信息，或者明知是编造的恐怖信息而故意传播，严重扰乱社会秩序的行为。

本罪是实害犯，要求造成严重扰乱了社会秩序的实害结果。

例1，甲对法院院长说：“如果不解决我的问题，我就给法院安装炸弹。”不构成本罪。

例2，甲对公安局局长谎称：“我已经在市政府安装了炸弹。”市政府紧急排查。甲构成本罪。

3.编造、故意传播虚假信息罪（第291条之一）

这是指，编造虚假的险情、疫情、灾情、警情，在信息网络或者其他媒体上传播，或者明知是上述虚假信息，故意在信息网络或者其他媒体上传播，严重扰乱社会秩序的行为。

【注意】这种虚假信息与虚假的恐怖信息是不同的。

\subsection*{五、考试作弊类犯罪（第 284 条之一)}

\subsubsection*{（一）组织考试作弊罪}

1.根据司法解释，\footnote{2019年9月4日最高人民法院、最高人民检察院《关于办理组织考试作等刑事案件适用法律若干问题的解释》。}本罪中的“考试”，仅限于全国人大及其常委会制定的法律所规定的考试。例如，高考、考研、法考、注册会计师资格考试、医师资格考试、建造师资格考试、中央和地方公务员录用考试等。（2016年第60题）

第一，上述考试涉及的特殊技能测试、面试等考试，属于“法律规定的国家考试”。

第二，本罪的考试，不包括国外的考试。例如，托福、雅思等。

2.既遂。组织考试作弊，在考试开始之前被查获，但已经非法获取考试试题、答案或者具有其他严重扰乱考试秩序情形的，应当认定为组织考试作弊罪\CJKunderwave{既遂}。

\subsubsection*{（二）非法出售、提供试题、答案罪}

1.这里的试题、答案要求是真实的，不能是虚假的，但不要求完全真实，只要求部分真实。试题不完整或者答案与标准答案不完全一致的，不影响非法出售、提供试题、答案罪的认定。

2.出售、提供的对象包括参加考试的人，也包括不参加考试的人。接收试题、答案的人是否使用了试题、答案，不影响出售、提供者构成本罪。

3.出售、提供的时间是考试前或考试中。

\subsubsection*{（三）代替考试罪}

1.行为主体是替考者和应考者。二者是对向犯性质的共犯关系。（2016年第60题）

[注意］一方构成犯罪，并不意味着另一方必然构成犯罪。例如，甲生病住院，无法考试。甲的父亲乙让丙代替甲参加考试。对此甲不知情。甲作为应考者不构成本罪，丙作为替考者构成本罪，乙教唆替考者，构成本罪的教唆犯。

2.为代替考试而使用虚假的身份证的，成立牵连犯，择一重罪论处。

% Source PDF physical page 369; see audit/source_coverage.json.
3.罪数。触犯本节罪名的同时，又触犯非法获取国家秘密罪、\CJKunderwave{非法利用信息网络罪}\CJKunderwave{的，想象竞合，择一重罪论处}。

\subsection*{六、计算机犯罪}

\subsubsection*{（一）非法侵入计算机信息系统罪（第285条第1款）}

本罪的计算机信息系统，仅包括受特别保护的计算机系统，也即国家事务、国防建设、尖端科技领域的计算机信息系统，不包括一般的计算机信息系统。例如，不包括一般公司企业的计算机信息系统。

\subsubsection*{（二）非法控制计算机信息系统罪（第285条第2款）}

1.行为对象：本罪的计算机信息系统，包括一切计算机信息系统。（1）智能手机系统属于计算机信息系统。（2）远程定位监控系统（如汽车租赁公司对车辆的定位监控），属于计算机信息系统。（2021年试题）

2.行为方式：非法控制。例如（制造“肉鸡”），张竣杰向目标服务器植入木马程序（后门程序）进行控制，未造成系统功能的破坏。张竣杰构成非法控制计算机信息系统罪。本案是2020年最高法院指导案例（第145号）。

[提示］代为出售这种控制权（肉鸡），构成掩饰、隐瞒犯罪所得罪。

\subsubsection*{（三）破坏计算机信息系统罪（第286条）}

1.行为对象：本罪的计算机信息系统，包括一切计算机信息系统。

2.行为类型一。本条第1款规定：干扰系统功能，造成系统不能正常运行。（2026年试题）

例1（劫持域名），李丙对甲网站实施域名劫持行为，修改甲网站的IP指向，使想访问甲网站的用户只能访问乙赌博网站，致使甲网站不能正常运行。李丙劫持甲网站的域名，干扰了甲网站的正常运行，构成破坏计算机信息系统罪。但是，该行为没有对甲网站本身进行控制，不构成非法控制计算机信息系统罪。（最高检察院指导案例第33号）

例2（锁定手机案），甲利用技术手段，远程锁定乙的苹果手机，导致乙的手机系统无法正常运行，乙无法正常使用手机，甲然后以“解锁”为条件，敲诈勒索乙的财物。甲实施这类行为共十起。甲构成破坏计算机信息系统罪，还构成敲诈勒索罪，属于牵连犯关系，择一重罪论处。（最高检察院指导案例第35号）

3.行为类型二。本条第2款规定：对数据删、改、增，造成严重后果。

问题：第2款行为类型要成立本罪，是否要求造成系统不能正常运行？对此有肯定说和否定说之争。

折中说（多数说）：看系统和数据是否关系到公共管理事务。如果是，并造成严重后果，即使没有造成系统不能正常运行，也能够成立本罪。

例1（删除违章信息案），交通管理员甲非法侵入交通违法信息管理系统，删除大量车辆违章信息，造成严重后果。甲构成本罪（最高法院研究室的复函意见）。其主要理由是，交通违法信息管理系统及其数据关系到公共管理事务，删改其中数据会造成严重后果。

例2（篡改志愿案），甲盗用高考考生乙的信息资料，登录高考招生系统，篡改乙的报考志愿，导致乙没有被所报考的高校录取，这属于造成严重后果。判例认为甲构成本

% Source PDF physical page 370; see audit/source_coverage.json.
罪。其主要理由是，高考招生系统及其数据关系到公共管理事务，删改其中数据会造成严重后果。（2024年试题）

例3（删差评案），李骏杰冒用某购物网站的许多买家的身份，骗取客服审核通过后，登录该网站内部评价系统，删改买家的差评347个。2017年最高检察院指导案例（第34号）认为，李骏杰构成破坏计算机信息系统罪。2024年，最高检察院宣布该指导案例失效，也即不认为李骏杰构成本罪。其可能的主要理由是，这种购物网站不属于关系公共管理事务的系统，删除差评不会给公共管理造成严重后果，并且没有造成系统不能正常运行。如果李骏杰以营利为目的，专门提供有偿删帖服务，则构成非法经营罪。

\subsubsection*{（四）非法获取计算机信息系统数据罪（第285条第2款）}

1.行为对象：一切计算机信息系统中的数据。

2.行为方式：（1）获取系统内已存的数据；（2）先在系统内生成数据，然后获取该数据。

例如（充值获取游戏币案），吴某发现某款网络游戏的充值系统存在漏洞，便使用技术手段为自已的游戏账号充值，获取大量游戏币，将部分游戏币出售牟利。吴某的行为属于先在系统内生成数据，然后获取该数据。吴某构成本罪。（《刑事审判参考》第1459号案例）（2020年试题）

[注意1]有些数据也属于虚拟财产（要求具有经济价值性和可管理性），因此本罪与盗窃罪不是对立排斥关系，可以想象竞合。

例如（盗窃域名案），张四毛欲窃取乙拥有的域名“www.8.cc”，利用技术手段将该域名转移到自已的账户，然后出售牟利。网络域名具有经济价值、可管理性，属于虚拟财产。张四毛构成盗窃罪。（最高检察院指导案例第37号）如果张四毛的技术手段同时构成非法获取计算机信息系统数据罪，则与盗窃罪想象竞合，择一重罪论处。同理，上述“充值获取游戏币案”中，吴某有可能构成盗窃罪，与非法获取计算机信息系统罪想象竞合，择一重罪论处。（2026年试题）

[注意2]代为出售这种数据，构成掩饰、隐瞒犯罪所得罪。

\subsubsection*{（五）非法利用信息网络罪（第287条之一）}

这是指，非法利用网络用于犯罪活动。

1.用于违法活动，不构成本罪（多数说）。

2.想象竞合。例如，在网上发布销售毒品的信息，同时触犯本罪和贩卖毒品罪，想象竞合，择一重罪论处。

\subsubsection*{（六）帮助信息网络犯罪活动罪（第287条之二）}

1.帮助方式，是提供技术支持、支付结算等帮助。

2.帮助对象，是信息网络犯罪活动。

3.成立本罪，需要遵守共犯从属性原理（多数说）。如果实行者不构成犯罪，则帮助者也不构成犯罪。例如，乙想实施网络诈骗，甲为乙提供网络接入服务，但乙后来又不想实施了。缺少乙的犯罪行为，甲的行为不构成本罪。（2022年试题）

4.想象竞合。例如，甲明知乙在网上开设赌场，仍为其提供网络接入服务。甲同时触犯本罪和开设赌场罪（帮助犯），想象竞合，择一重罪论处。（2016年第54题、2023年试题，2025年主观题）

% Source PDF physical page 371; see audit/source_coverage.json.
\section{妨害司法罪}

\subsection*{一、伪证罪}

第305条在\CJKunderwave{刑事诉讼}中，证人、鉴定人、记录人、翻译人对与案件有重要关系的情节，故意作虚假证明、鉴定、记录、翻译，意图陷害他人或者隐匿罪证的，处三年以下有期徒刑或者拘役；····

1.行为主体：刑事诉讼中的证人、鉴定人、记录人、翻译人。

本罪的“证人”应作扩大解释，既包括“证人证言”中的证人，也包括“被害人陈述”中的被害人。被害人在刑诉中故意作虚假陈述，构成本罪。

注意：本罪的证人，（1）既包括知道案件情况的人，也包括不知道案件情况的人。例如，甲不在现场，却故意作伪证，构成伪证罪。（2）知道案件情况的人，既包括知道案发情况，也包括知道“侦查案件”的情况。例如，侦查人员在看守所刑讯逼供嫌疑人，获取口供。看守所的监管人员甲看到这一幕。甲在法庭上作伪证，谎称侦查人员没有刑讯逼供，导致口供没有被排除。甲构成伪证罪。

2.发生领域：刑事诉讼中。发生时间：是否启动立案侦查一诉讼程序终结。

[注意1]公安机关在决定是否立案，作伤情鉴定时，如果鉴定人作虚假鉴定，构成本罪。

[注意2］终审后，启动审判监督程序。在该程序中，作伪证，构成本罪。

3.伪证的内容，是与案件有重要关系的情节。该情节包括定罪情节和量刑情节。例如，证人甲就被告人是否自首、立功作伪证，构成伪证罪。

4.主观是故意。（1）没有故意作伪证，陈述内容却与客观事实不符，不构成伪证罪。

（2）故意作伪证，却与客观事实相符，不构成伪证罪。（3）司法人员甲在办理乙涉嫌盗窃案时，要求丙作伪证，丙只好照办。丙有故意，但是缺乏期待可能性，不以伪证罪论处。甲构成伪证罪的教唆犯（伪证罪的教唆犯被正犯化为妨害作证罪，定妨害作证罪）和徇私枉法罪的实行犯，想象竞合，择一重罪论处。

5.既遂标准：作完伪证。成立既遂，不要求使法官相信、作出错误判决。（2023年试题）。例如，证人甲在第一次询问程序中作完伪证；在第二次询问程序中，又修正第一次的伪证。第一次的伪证罪依然既遂。

6.伪证罪与诬告陷害罪的关系。

（1）诬陷到底的情形。例如，在立案前，甲向司法机关虚假告发，谎称“乙强奸了丙”。甲构成诬告陷害罪。司法机关立案后，甲又向司法机关作伪证，谎称“自己是目击证人，看到乙强奸丙”。甲构成伪证罪。多数说认为，前后两个行为构成两个罪，应并罚。

（2）落井下石的情形。例1，丙因盗窃罪被抓。甲向司法机关诬告丙还犯有强奸罪。甲构成诬告陷害罪，不构成伪证罪。例2，司法机关在审理丙涉嫌盗窃案时，证人乙对丙的盗窃案作伪证，构成伪证罪；乙又谎称丙还犯有强奸罪，构成诬告陷害罪。两罪并罚。

\subsection*{二、妨害作证罪}

第307条第1款以暴力、威胁、贿买等方法阻止证人作证或者指使他人作伪证的，

% Source PDF physical page 372; see audit/source_coverage.json.
处三年以下有期徒刑或者拘役；·

1.行为主体：一般主体。司法工作人员犯本罪，从重处罚。

2.发生领域：多数说认为包括刑事诉讼、民事诉讼、行政诉讼。\footnote{少数说认为只包括刑事诉讼。}

3.妨害作证：第一，阻止证人作证。第二，指使他人作伪证。（2007年第64题、2023年试题）

［注意］这里的“作证”应扩大解释，包括：证人证言，被害人陈述，记录人，翻译人，鉴定意见，勘验、检查笔录。

例1，甲暴力阻止被害人作陈述，阻止鉴定人作鉴定，阻止勘验、检查人员作勘验、检查。甲构成本罪。

例2，甲用威胁手段，指使被害人作虚假陈述，指使鉴定人作虚假鉴定，指使勘验、检查人员作虚假的勘验、检查笔录。甲构成本罪。

例3，甲用威胁手段，指使不知道案情的乙作伪证。甲构成本罪。

4.妨害手段：以暴力、胁迫、贿买等方法。这里的“等方法”包括教唆、指使。（本犯人的手段有专门分析，见下文）

例如，甲、乙是好友，乙的重大贪污罪行被丙发现。甲是丙的上司，为防止丙作证，将丙派往国外工作。甲的行为虽然防止了丙作证，但不属于以暴力、胁迫、贿买、教唆等方法阻止丙作证，所以不构成本罪。（2014年第61题）

5.既遂标准：客观上阻止了证人作证；使他人作出了伪证。

6.本罪是将伪证罪的教唆犯给予正犯化了。例如，甲教唆证人乙作伪证。乙照办。乙构成伪证罪，甲构成伪证罪的教唆犯。但是，由于有了独立的妨害作证罪，所以，对甲不定伪证罪的教唆犯，而定妨害作证罪。

7.犯罪嫌疑人（本犯人）的期待可能性问题。例如，甲乙共同盗窃，被抓。

（1）甲作虚假供述的，不构成伪证罪，因为法律不期待其能够作出真实供述。

（2）甲指使丙作伪证，丙照办。丙构成伪证罪。

第一，甲是否构成伪证罪的教唆犯？不构成。理由是，第一，共犯从属性原理仅表明，实行者构成犯罪，是教唆者构成犯罪的必要条件（前提条件），而不是充分条件。第二，甲自己作虚假供述，由于缺乏期待可能性而不构成伪证罪的实行犯，那么甲的教唆行为比实行行为轻，就更不应该构成伪证罪的教唆犯。\footnote{少数说认为构成。理由是，根据共犯从属性原理，既然丙构成伪证罪的实行犯，那么甲便构成伪证罪的教唆犯。}

第二，甲是否构成妨害作证罪？\CJKunderwave{甲如果采取普通的教唆行为，指使丙作伪证，或者阻}\CJKunderwave{止丙作证，则不构成妨害作证罪。甲如果采取暴力、威胁、贿买这三种手段，指使丙作伪}\CJKunderwave{证，或者阻止丙作证，则构成妨害作证罪}。这是因为，还是能够期待甲不要采取这三种严重手段。\footnote{少数说认为，既然甲不构成伪证罪的教唆犯，那么也不构成妨害作证罪。}

（3）甲指使共犯人乙作虚假供述。乙照办。乙不构成伪证罪。甲不构成伪证罪的教唆犯。问题是，甲是否构成妨害作证罪？甲如果采取普通的教唆行为，指使乙作虚假供述，或者阻止乙作供述，则不构成妨害作证罪。甲如果采取暴力、威胁、贿买这三种手段，指

% Source PDF physical page 373; see audit/source_coverage.json.
使乙作虚假供述，或者阻止乙作供述，则构成妨害作证罪。这是因为，还是能够期待甲不要采取这三种严重手段。\footnote{少数说认为，既然甲不构成伪证罪的教唆犯，那么也不构成妨害作证罪。}

8.罪数。例1，甲用暴力手段阻止勘验人、检查人（司法工作人员）进行勘验、检查。甲不仅构成妨害作证罪，还构成妨害公务罪，想象竞合，择一重罪论处。

例2，甲用贿买的手段，给勘验人乙（司法工作人员）送钱，指使乙不要勘验现场，以便让甲的儿子的杀人证据灭失。乙照办。乙构成受贿罪和徇私枉法罪，根据第399条第4款，择一重罪论处。甲构成行贿罪，徇私枉法罪的教唆犯，妨害作证罪，想象竞合，择一重罪论处。

\subsection*{三、帮助毁灭、伪造证据罪}

第307条第2款帮助当事人毁灭、伪造证据，情节严重的，处三年以下有期徒刑或者拘役。

1.行为主体：一般主体。司法工作人员犯本罪，从重处罚。

2.发生领域：多数说认为，包括刑事诉讼、民事诉讼、行政诉讼。发生时间：可以诉讼前，可以诉讼后，没有时间限制。

例1，甲故意砸坏乙的轿车，后感到害怕，让丙修理，修理费是1万元。乙发现后指控甲构成故意毁坏财物罪。警方在决定是否立案时，要评估毁坏财物的价值是否达到立案标准（如5千元），请丙提供修理费证明。甲用普通教唆手段，指使丙出具修理费为4千元的虚假证明。丙照办。警方因此没有立案。丙构成帮助伪造证据罪、包庇罪。甲是本犯人，用普通教唆手段，不构成妨害作证罪，因为不具有期待可能性。

例2，甲被终审后判有罪。乙为了申诉，帮助甲伪造证据。乙构成帮助伪造证据罪。

例3，甲杀人后潜逃。在潜逃期间，乙帮助甲毁灭证据，构成帮助毁灭证据罪。

3.帮助对象：当事人。

（1）这里的当事人不包括自己。自己为自己，不算帮助。自己犯罪，为自己毁灭、伪造证据，不具有期待可能性，不构成犯罪。

问题是，甲教唆乙为自己（甲）毁灭证据，乙照办。乙构成帮助毁灭证据罪。甲是否构成帮助毁灭证据罪的教唆犯？不构成。理由是，第一，共犯从属性原理仅表明，实行者构成犯罪，是教唆者构成犯罪的必要条件（前提条件），而不是充分条件。第二，甲自己毁灭证据，由于缺乏期待可能性而不构成本罪的实行犯，那么甲的教唆行为比实行行为轻，就更不应该构成本罪的教唆犯。\footnote{少数说认为构成。理由是，根据共犯从属性原理，既然乙构成本罪的实行犯，那么甲便构成本罪的教唆犯。}

（2）这里的当事人是否包括共犯人？行为人如果专门为了共犯人，则构成本罪；如果专门为自己，或既为自己，也为共犯人，则不构成本罪。例如（2014年第61题），甲、乙共同盗窃了丙的财物。为防止公安人员提取指纹，甲在丙报案前擦掉了两人留在现场的指纹。甲不构成本罪。\footnote{少数说认为，上述情形中，行为人都不构成本罪。}

4.“证据”范围。第一，包括实物证据和实体化的言词证据，例如，作为视听资料的证人证言、被害人陈述、嫌疑人供述。第二，不包括人。例如，甲将证人和被害人藏匿起

% Source PDF physical page 374; see audit/source_coverage.json.
来，阻止其作证。甲不构成本罪，但构成妨害作证罪。

5.帮助类型。“毁灭证据”包括隐匿证据。“伪造证据”包括变造证据。这都是扩大解释。

（1）在刑事诉讼中，即使经过当事人（犯罪嫌疑人、被告人）同意，帮助其毁灭无罪证据、伪造不利证据，由于妨害了刑事司法客观公正性，也构成帮助毁灭、伪造证据罪。（2014年第61题）

（2）在民事诉讼、行政诉讼中，经当事人同意，帮助其毁灭有利证据、伪造不利证据，因为民事、行政诉讼尊重当事人的处分权，所以行为人不构成帮助毁灭、伪造证据罪。

6.帮助行为。“帮助”应作扩大解释，包括狭义的帮助行为，还包括实行行为、教唆行为。这个考点很重要。

[总结］甲杀了人，乙是帮助者：（2009年第63题，2014年第61题，2015年第20题）

\InsertDiagram{figures/manual_table_p374_1.tex}

7.既遂标准：毁灭掉证据、伪造好证据。

\subsection*{四、虚假诉讼罪（第307条之一）}

\InsertDiagram{figures/manual_table_p374_2.tex}

% Source PDF physical page 375; see audit/source_coverage.json.
\subsection*{五、窝藏、包庇罪（第310条）}

窝藏罪，是指明知是犯罪的人，帮助其逃匿的行为。包庇罪，是指明知是犯罪的人，作假证明包庇的行为。

\subsubsection*{（一）窝藏、包庇的对象：犯罪的人}

这里的“犯罪的人”，包括犯罪嫌疑人、被告人。也即，这里的“犯罪的人”，不能从无罪推定出发认为判决有罪之前的人都是无罪的人，都不是这里的“犯罪的人”。

1.犯罪嫌疑人归案后被宣告无罪的，应当依照法定程序宣告窝藏、包庇行为人无罪。\footnote{2021年8月11日最高人民法院、最高人民检察院《关于办理窝藏、包庇刑事案件适用法律若干问题的解释》。}例如，甲未杀人，因冤枉而被列为犯罪嫌疑人。乙将甲藏起来。后甲乙均被抓捕。甲被宣告无罪。乙也不构成窝藏罪。

2.刑法第362条规定：“旅馆业、饮食服务业、文化娱乐业、出租汽车业等单位的人员，在公安机关查处卖淫、活动时，为违法犯罪分子通风报信，情节严重的，定窝藏、包庇罪。”\footnote{提示：2017年7月25日最高人民法院、最高人民检察院《关于办理组织、强迫、引诱、容留、介绍卖淫刑事案件适用法律若干问题的解释》第14条认为，对这种情形以包庇罪论处。然而，通风报信属于帮助逃匿，应定窝藏罪。所以，该司法解释只能被视为特殊规定。}

在此，如果行为人通风报信的是卖淫、等违法行为，则该条是法律拟制。因为卖淫、者是违法分子，不是犯罪分子，而窝藏、包庇罪要求的行为对象是犯罪分子。如果行为人通风报信的是强迫卖淫、宿幼女等犯罪行为，则该条是注意规定。因此，该条既包含法律拟制的情形，也包含注意规定的情形。（2009年第62题）

\subsubsection*{（二）窝藏罪}

这是指，帮助犯罪人逃匿的行为。这里的“帮助”应作扩大解释，包括狭义的帮助行为，还包括实行行为、教唆行为。

1.实行行为。例如，将犯罪人藏在自已家里；开车将犯罪人运到外地。

2.狭义的帮助行为。例如，犯罪人要逃往外地，为其提供车辆、钱财、虚假身份证、化妆用具，为其通报侦查或追捕消息。（2012年第19题）

[注意1]帮助行为要对逃匿行为起到实质的、直接的帮助效果。也即，给司法机关抓捕犯罪人直接增加难度。

例1，向犯罪人提供一个关公护身符，保佑其过五关，斩六将，不成立窝藏罪。向犯罪人提供刀具。多数说认为，该行为不会给司法机关直接增加难度，不构成窝藏罪。

例2，犯罪人甲的妻子陪同甲逃到外地，为甲洗衣做饭，照顾甲的起居，不构成窝藏罪。（2023年试题）

例3，受已经逃匿于外地的犯罪人之托，向犯罪人妻子提供金钱，使犯罪人安心逃匿的行为，不成立窝藏罪。（2023年试题）

[注意2］帮助行为可以违背犯罪人意愿。例如，犯罪人不想逃匿，甲将犯罪人弄晕，然后将其运至外地，构成窝藏罪。（2023年试题）

3.教唆行为。例如，甲劝诱犯罪人逃往外地，犯罪人答应照办。甲构成窝藏罪。提

% Source PDF physical page 376; see audit/source_coverage.json.
示：第一，如果犯罪人拒绝了逃匿，这种劝诱行为不值得定罪。第二，犯罪人意欲自首，行为人仅劝诱其不自首，但没劝诱逃匿的，不成立窝藏罪。（2023年试题）

\subsubsection*{（三）包庇罪}

这是指，向司法机关作假证明，帮助犯罪人逃避刑事追究或者获得从宽处罚。

1.包庇的目的包括两种：一是定罪方面，帮助犯罪人逃避刑事追究；二是量刑方面，帮助犯罪人获得从宽处罚。

2.总结“知情不报”：

（1）知情不报，是指明知发生犯罪事实或知道犯罪人去向，但是不举报。知情不报不构成窝藏、包庇罪。在司法机关调查取证时，仍知情不报，也不构成窝藏、包庇罪。在司法机关调查取证时，故意提供虚假证据，帮助犯罪人，则构成包庇罪或伪证罪。

（2）保证人在犯罪人取保候审期间，明知犯罪人的藏匿地点、联系方式，但拒绝向司法机关提供的，对保证人以窝藏罪论处。\footnote{2021年8月11日最高人民法院、最高人民检察院《关于办理窝藏、包庇刑事案件适用法律若干问题的解释》。}

（3）明知他人有间谍犯罪、恐怖犯罪、极端主义犯罪行为，在司法机关向其调查时，拒绝提供，情节严重的，构成拒绝提供间谍犯罪、恐怖主义犯罪、极端主义犯罪证据罪（第311条）。作假证明包庇的，定包庇罪。（2012年第19题）

3.包庇罪与伪证罪，可以产生想象竞合。例如，向公安司法机关作伪证帮助犯罪人逃避刑事追究，就同时触犯了包庇罪和伪证罪，属于想象竞合，择一重罪论处。

4.包庇罪与帮助毁灭、伪造证据罪，可以产生想象竞合。例如，甲在移动公司办业务时，与营业员发生口角，甲一气之下砸毁办公设备，维修费用是1万元。移动公司报警，警察立案。甲找到乙，让乙帮忙。乙找到移动公司负责人丙，丙答应将维修费用单改为4千元。由于数额不够5千元，警察撤案。甲构成故意毁坏财物罪。乙丙构成包庇罪和帮助伪造证据罪，想象竞合，择一重罪论处。

5.特殊类型的包庇犯罪。（2006年第63题）

包庇罪与特殊类型的包庇犯罪是法条竞合关系，优先适用这些特殊罪名：（1）包庇黑社会性质组织罪（第294条第3款，这是指国家机关工作人员包庇黑社会性质组织）。（2）包庇毒品犯罪分子罪（第349条）。

\subsubsection*{（四）主观要件}

主观上要求明知是犯罪的人。注意：行为人将犯罪人所犯之罪误认为其他犯罪的，不影响“明知”的认定。例如，甲以为乙犯了抢劫罪而窝藏乙，实际上乙犯了强奸罪。甲依然构成窝藏罪。\footnote{根据司法解释规定，虽然为犯罪的人提供隐藏处所、财物，但不是出于帮助犯罪人逃匿的目的，不以窝藏罪论处。对未履行法定报告义务的行为人，给予行政处罚。该规定主要是针对酒店而言的。例如，酒店明知客人可能是逃犯，由于不便调查，便办理了入住手续。酒店不构成窝藏罪。}

\subsubsection*{（五）认定问题}

1.窝藏、包庇罪属于抽象的危险犯。因此，即使警方知道犯罪人被行为人隐藏在何处，即使警方明知行为人提供的是虚假证明，也不妨碍窝藏、包庇罪的成立。

2.期待可能性。（1）犯罪人窝藏自己、包庇自己，不具有期待可能性，不构成窝藏、

% Source PDF physical page 377; see audit/source_coverage.json.
包庇罪。（2）犯罪人甲教唆乙窝藏、包庇甲，乙照办。乙构成窝藏、包庇罪。甲不构成教唆犯，因为不具有期待可能性。

3.共犯人之间窝藏、包庇，根据司法解释，不构成窝藏、包庇罪，因为不具有期待可能性。

4.罪数：“一条龙服务”。为帮助同一个犯罪的人逃避刑事处罚，实施窝藏、包庇行为，又实施洗钱行为，或者掩饰、隐瞒犯罪所得及其收益行为，或者帮助毁灭证据行为，或者伪证行为的，依照重罪吸收轻罪处理，不实行数罪并罚。

\subsection*{六、掩饰、隐瞒犯罪所得、犯罪所得收益罪（第312条）}

本罪是指，明知是犯罪所得及其收益，而予以窝藏、转移、代为销售、收购或者以其他方法掩饰、隐瞒的行为。本罪的保护法益是司法机关的追赃活动。

\subsubsection*{（一）行为主体}

1.行为主体仅限于本犯以外的人。这里的本犯包括上游犯罪的实行犯、教唆犯、帮助犯。这些本犯实施赃物犯罪，不构成本罪，因为不具有期待可能性。换言之，事前有通谋，事后掩饰、隐瞒犯罪所得的，不构成本罪，而构成上游犯罪的共犯。（2011年第17题）

例1，甲盗窃轿车，然后藏在家里。对甲只以盗窃罪论处，不再定本罪。

例2，甲教唆乙盗窃轿车，乙窃得后交给甲，甲窝藏了轿车。对甲只以盗窃罪（教唆犯）论处，不再定本罪。

2.行为主体包括单位。盗用单位名义，实施本罪行为，违法所得归个人所有的，由于没有体现单位意志，没有为单位谋利，因此按个人犯罪处理。

\subsubsection*{（二）行为对象}

行为对象：犯罪所得及其产生的收益。

1.犯罪“所得”的范围。

犯罪所得，是指通过犯罪行为直接获得的赃物，包括财物和财产性利益。犯罪所得的收益，是指利用犯罪所得获得的财产性收益。例如，将受贿款存入银行后所获得的利息，属于犯罪所得产生的收益。

（1）虚拟财产能够成为赃物。

例1，甲盗窃虚拟游戏币，乙帮助窝藏的，构成本罪。

例2，甲实施非法获取计算机信息系统数据罪，获取的数据视为赃物。乙代为销售的，构成本罪。

例3，甲实施非法控制计算机信息系统罪，获取控制权。乙代为出售该控制权的，构成本罪。例2、例3是司法解释的规定。

（2）违禁品能够成为赃物。例如，甲盗窃了大量管制刀具，乙知道真相而予以收购的，构成本罪。

（3）人的身体不是赃物。例如，甲收买了被拐卖的妇女、儿童后，得知警察要来解救，请求乙窝藏，乙照办，将妇女、儿童关押在小黑屋里。乙不属于窝藏“犯罪所得”。乙构成非法拘禁罪，对儿童还构成拐骗儿童罪（与非法拘禁罪想象竞合）。

（4）犯罪工具不是赃物。例如，甲骑摩托车抢夺了乙的项链。甲向丙告知真相，让丙帮助将摩托车和项链藏起来。项链是甲的犯罪所得，丙对项链构成掩饰、隐瞒犯罪所得

% Source PDF physical page 378; see audit/source_coverage.json.
罪。摩托车是甲的犯罪工具，此时属于犯罪证据，丙对摩托车构成帮助毁灭证据罪。

2.“犯罪”所得的范围，也即上游犯罪的范围。

（1）包括财产犯罪、经济犯罪的犯罪所得。例如，通过盗窃罪获取的文物、油气等。

（2）包括其他能够获取财物的犯罪。例如，通过赌博罪赢得的钱财，通过受贿罪获取的贿赂，通过非法狩猎罪获取的动物，通过非法捕捞水产品罪获取的水产品。

（3）包括收购赃物型掩饰、隐瞒犯罪所得罪。例如，甲盗窃到原油，乙知道真相而购买，乙买到后交给丙代为销售，丙知道真相而照办。乙构成收购型掩饰、隐瞒犯罪所得罪，买到后，原油便成为乙的掩饰、隐瞒犯罪所得罪的犯罪所得。丙代为销售，构成掩饰、隐瞒犯罪所得罪。假如乙买到后，自已窝藏或销售，则这种窝藏或销售行为不再构成掩饰、隐瞒犯罪所得罪，因为不具有期待可能性，属于不可罚的事后行为。

（4）包括制造违法事实阶段的犯罪。例如，15周岁的甲盗窃到摩托车，16周岁的乙代为销售。乙构成本罪。

（5）不包括挪用公款罪，因为该罪行为人是想将公款还回去的，所以，所挪用的公款不算犯罪“所得”。犯罪所得，要求据为己有。

[提示］法考真题答案认为，挪用公款罪能够成为洗钱罪的上游犯罪。

\subsubsection*{（三）行为方式}

本罪的行为方式有窝藏、转移、代为销售、收购，或者其他掩饰、隐瞒的行为。共同特征是，妨害司法机关的追赃活动。

例1，甲职务侵占了100万元单位钱款，交给知道真相的儿子乙。乙用来购买房屋或者消费掉，乙构成本罪。

例2，甲盗窃到汽车。乙帮助拆解、组装、更改车身颜色、修改发动机号、提供机动车来历凭证，构成本罪。（2011年第17题）

[巩固练习]题1，甲公司收购了他人盗窃来的原油，甲公司的技术人员乙知道真相，对原油进行质量检测。甲公司构成本罪（本罪可由单位构成）。乙不构成本罪，因为乙的行为属于中立行为，没有妨害司法机关的追赃活动。

题2，非国家工作人员甲受贿获取一套房屋。乙装修公司知道真相，对该房屋进行装修。乙公司不构成本罪，因为其行为属于中立行为，没有妨害司法机关的追赃活动。

题3（帮助行为），甲为了窝藏、转移自已所盗窃的大型赃物，需要卡车，乙知道真相仍将卡车提供给甲，使甲顺利窝藏、转移了赃物。甲乙在制造违法事实阶段构成赃物犯罪的共同犯罪，甲是实行犯，乙是帮助犯；在排除责任事由阶段，甲由于不具有期待可能性而最终不构成本罪，乙最终构成本罪（帮助犯）。

\subsubsection*{（四）参与时间}

本罪的成立时间是上游犯罪既遂后。若上游犯罪尚未既遂，行为人参与进来，则构成上游犯罪的承继的共同犯罪。

例1，甲抢劫乙的钱包，将乙打倒在地，丙看到后参与进来，帮助甲捡起钱包。丙构成抢劫罪的承继的共犯，而非本罪。

例2，甲通过电信诈骗获得一万元，告知乙真相，让乙帮助存在乙的银行卡里。乙照办。乙构成本罪。此后，甲想继续实施电信诈骗，让乙提供银行卡，乙知道真相仍提供，

% Source PDF physical page 379; see audit/source_coverage.json.
甲让受骗人将钱款打到乙的卡里。乙取出给甲。乙构成诈骗罪的共同犯罪和帮助信息网络犯罪活动罪，想象竞合。\footnote{2025年7月22日最高人民法院、最高人民检察院、公安部《关于办理帮助信息网络犯罪活动等刑事案件有关问题的意见》。}

例3，父亲持信用卡的主卡，女儿持副卡消费。女儿消费后，由父亲利用职务侵占的所得款归还，女儿知情。女儿不构成本罪，因为女儿透支消费的款项源于发卡银行，父亲是向发卡银行偿还。如果父亲将职务侵占的钱款直接打到女儿卡里，女儿消费，则女儿构成本罪。

\subsubsection*{（五）主观故意}

对象错误：行为人误以为是A种上游犯罪所得，实际是B种上游犯罪所得，予以掩饰、隐瞒，构成本罪。例如，乙偷了丙的车，甲以为是乙抢劫来的车，帮乙卖了出去。甲构成本罪。（2024年试题）

\subsubsection*{（六）上游犯罪的要求}

1.认定本罪，以上游犯罪事实存在为前提。上游犯罪事实查证属实，但行为人尚未到案，或者因行为人死亡、无责任能力等原因不追究刑事责任的，不影响本罪的认定。\footnote{2025年8月26日最高人民法院、最高人民检察院《关于办理掩饰、隐瞒犯罪所得、犯罪所得收益刑事案件适用法律若干问题的解释》。}

2.上游犯罪的数额与本罪的数额

（1）上游犯罪没达到定罪数额，为之实施掩饰、隐瞒行为，不构成本罪。这是因为，成立本罪，要求上游行为构成犯罪。例如，甲实施职务侵占，获得1万元财物。职务侵占罪的定罪数额是3万元。甲不构成职务侵占罪。因此，乙帮助窝藏该1万元财物，不构成本罪。

（2）上游犯罪达到定罪数额，为之实施掩饰、隐瞒行为（销赃行为），获得价值数额大小，不影响本罪的成立。这是因为，本罪不是财产犯罪，而是妨害司法秩序的犯罪。

例1，甲盗窃到一辆摩托车（价值1万元）。乙代为销售，销售水平差，仅以100元售出（气得甲吐血），买家知道是赃物。乙构成本罪，因为妨害了司法机关的追赃活动。（旧司法解释要求销赃数额达到3千才定罪。该规定已被删除）

例2，甲抢劫到500元财物。乙代为销售，销售水平差，仅以100元售出，买家知道是赃物。甲构成抢劫罪，因此，乙构成本罪。

（3）本罪的数额，应当以实施掩饰、隐瞒行为时为准。收购或者代为销售财物的价格高于其实际价值的，以收购或者代为销售的价格计算。\footnote{2025年8月26日最高人民法院、最高人民检察院《关于办理掩饰、隐瞒犯罪所得、犯罪所得收益刑事案件适用法律若干问题的解释》。}

\subsubsection*{（七）销赃的罪数总结}

\InsertDiagram{figures/manual_table_p379_1.tex}

% Source PDF physical page 380; see audit/source_coverage.json.
\subsubsection*{（八）“黑吃黑”的问题}

成立本罪，要求不违背上游犯罪人（本犯人）的意愿（多数说）。若对本犯人的犯罪所得及其收益，实施盗窃、抢劫、诈骗、抢夺行为，则构成相应犯罪。也即，犯罪人对犯罪所得及其收益的非法占有能够对抗一般人。\footnote{2025年8月26日最高人民法院、最高人民检察院《关于办理掩饰、隐瞒犯罪所得、犯罪所得收益刑事案件适用法律若干问题的解释》。}若实施侵占行为，则参见财产犯罪一讲侵占罪一节。（2025年试题）

\subsection*{七、脱逃罪（第316条）}

1.行为主体：依法被关押的罪犯、犯罪嫌疑人、被告人。一是指已被判处拘役以上刑罚的正在监狱服刑的已决犯。二是指已经被刑事拘留、逮捕而未审判的未决犯，即犯罪嫌疑人、被告人。

（1）被行政拘留、司法拘留的人不是罪犯，不是本罪主体。

（2）被采取\CJKunderwave{拘传、取保候审、监视居住}等强制措施的犯罪嫌疑人、被告人，因为不在关押状态，不是本罪主体。

（3）被判处管制、有期徒刑的缓刑、被假释的罪犯，因为不在关押状态，不是本罪主体。这些人如果违反相关规定，以相关规定处理。例如，撤销缓刑、假释。

（4）被公民扭送的现行犯、通缉犯、越狱逃跑犯、正在被追捕的人，挣脱公民扭送的，不构成本罪。

(5）被监察机关采取留置措施的犯罪嫌疑人，是本罪主体。（2019年试题）

（6）司法解释规定：对被错判徒刑的在服刑期间脱逃的行为，可不以脱逃罪论处；但在脱逃期间犯罪的，应定罪判刑；对被错判已服刑日期与后来犯罪所判处的刑期不宜折抵，可在量刑时酌情考虑从轻或减轻处罚。\footnote{1983年8月31日最高人民法院研究室《关于因错判在服刑期脱逃后确有犯罪其错判服刑期限可否与后判刑期折抵问题的电话答复》。}

2.认定问题。

（1）脱逃地点：一般在监狱、看守所等关押场所。行为人在被转监（从一个监狱转到另一个监狱）、送往法院审判途中脱逃的，也构成本罪。

（2）行为人受到奖励，节假日获准回家，故意不在规定时间返回监狱，逃往外地，构成本罪。如果是因为特殊原因，例如，因父母重病而迟滞不归，特殊原因消除后返回，不构成本罪。

（3）既遂标准：摆脱监管人员的实力控制。逃出关押场所，但明显处于被追捕过程中，没有摆脱被追捕的，仍属于未遂。例如，囚犯甲翻墙逃出监狱，但是当即被发现，甲在前面逃跑，狱警在后面紧追，此时甲虽逃出监狱但仍属未遂。如果甲逃出监狱并摆脱追捕状态，大隐于市，狱警不知甲身藏何处，到处张贴通缉令。甲构成脱逃罪既遂。

（4）共犯问题：甲乙商议共同脱逃，只要一个人既遂，根据“部分实行、全部负责”原则，未脱逃成功的或中止脱逃的行为人，也认定为既遂。（2018年试题）

\subsection*{八、拒不执行判决、裁定罪（第313条）}

1.本罪是指，负有执行义务的人对法院的判决、裁定有能力执行而拒不执行，情节严重。

% Source PDF physical page 381; see audit/source_coverage.json.
2.本罪是不作为犯，应为+能为+不为+情节严重。

3.隐藏、转移财产的时间点问题。行为人为逃避执行义务，在诉讼开始后、裁判生效前，实施隐藏、转移财产的行为，在判决、裁定生效后，拒不执行的，构成本罪。\footnote{2024年12月1日最高人民法院、最高人民检察院《关于办理拒不执行判决、裁定刑事案件适用法律若干问题的解释》第6条；参见张明楷：《拒不执行判决、裁定罪的基本问题》，载《中国刑事法杂志》2025年第1期。}（2025年试题）

\section{妨害文物管理罪}

\subsection*{一、倒卖文物罪（第326条）}

1.行为对象：国家禁止经营的文物，而不是所有文物。

2.倒卖，根据司法解释，\footnote{2016年1月1日最高人民法院、最高人民检察院《关于办理妨害文物管理等刑事案件适用法律若干问题的解释》。}是指以牟利为目的，出售或者为出售而收购、运输、储存的行为。

（1）“牟利目的”只要求具有，不要求实现。

（2）倒卖不要求“买进+卖出”。这意味着：

第一，出售自己收藏的国家禁止经营的文物，构成倒卖文物罪。

第二，盗窃国家禁止经营的文物，然后出售，构成盗窃罪和倒卖文物罪，并罚。

\subsection*{二、盗掘古墓葬罪（第328条）}

1.盗掘，集盗窃与损毁于一体，是指未经国家文物主管部门批准，私自挖掘古文化遗址、古墓葬。（2025年试题）

2.本罪与盗窃罪的关系。盗掘行为本身包含了盗窃行为，所以盗掘过程中同时盗窃所掘取的文物，不再定盗窃罪，而是作为法定刑升格条件（2010年第63题）。此后，出售该文物，又构成倒卖文物罪，与盗掘古墓葬罪并罚。

3.本罪与故意损毁文物罪的关系。盗掘行为本身包含了损毁行为，所以盗掘过程中损毁文物，不再定故意损毁文物罪，而是作为法定刑升格条件。

\section{妨害国（边）境管理罪}

组织他人偷越国（边）境罪（第318条）和运送他人偷越国（边）境罪（第321条）

1.既遂标准：越过国（边）境线。（2022年试题）

2.罪数。（1）结合犯：本罪+妨害公务罪=本罪（加重处罚）。（2）结合犯：本罪+非法拘禁罪=本罪（加重处罚）。（3）本罪+其他犯罪，并罚。

% Source PDF physical page 382; see audit/source_coverage.json.
\section{危害公共卫生罪}

[问题]某乡村教师甲颇懂医术，常为自已治病。某日村民乙请求甲为其母治病，甲答应。由于甲的治疗手段不当导致病人死亡。甲是否构成非法行医罪？（颇懂医术案）

\subsection*{一、医疗事故罪}

第335条医务人员由于严重不负责任，造成就诊人死亡或者严重损害就诊人身体健康的，处三年以下有期徒刑或者拘役。

1.行为主体是特殊主体：医务人员。

2.行为方式既可以是作为，也可以是不作为。作为的方式，例如，因疏忽大意误将刺激性针剂当作麻醉针剂给病人注射，导致病人死亡。不作为的方式，例如，应当及时检查病情因为睡着了没有检查，导致病人死亡。

3.主观是过失。这里的过失属于一种业务过失，往往是重大过失。

4.成立本罪要求严重后果：造成就诊人死亡或者严重损害就诊人身体健康。

\subsection*{二、非法行医罪（第336条）}

1.行为主体：未取得医生执业资格的人。

（1）具有医生执业资格的人，不可能成为本罪的实行犯，但可以成为帮助犯或者教唆犯。例如，具有医生执业资格的甲教唆没有该资格的乙非法行医，同时为乙提供帮助。乙构成非法行医罪，甲构成非法行医罪的教唆犯和帮助犯。

（2）具有医生执业资格的人，在私人诊所或其他医院兼职，不构成本罪，造成就诊人员重伤、死亡的，可构成医疗事故罪。

2.行为表现：非法行医。

（1）本罪是典型的职业犯，要求有反复、持续实施的意思。例如，甲是药剂师，没有医生执业资格，但在药店长期坐堂治病，导致某个病人死亡。

[注意]不具有医生执业资格的人，没有反复、持续实施的意思，偶然为他人治病，不成立非法行医罪。例如，本节开头“颇懂医术案”中，甲不构成非法行医罪，而构成过失致人死亡罪。

（2）即使病人承诺接受治疗，行为人非法行医也可构成犯罪。

（3）非法行医过失致人死亡，构成结果加重犯，只定非法行医罪，加重处罚，不再定过失致人死亡罪。（2013年第18题）

\section{环境犯罪}

[问题]甲在盗伐林木时，被护林员乙发现，乙抓捕甲，甲为了抗拒抓捕，将乙打成重伤。对甲如何处理？（打护林员案）

% Source PDF physical page 383; see audit/source_coverage.json.
\subsection*{一、针对林木的犯罪（第345条）}

区分原理：盗窃罪=A（非法占有他人林木，侵犯财产权），盗伐林木罪=A+B（破坏生态环境），滥伐林木罪=B（破坏生态环境，没有非法占有目的）。

\subsubsection*{（一）盗伐林木罪}

1.行为公式：盗伐林木罪=A（盗窃罪）+B（破坏生态环境）。

根据司法解释，盗伐林木罪，是指带着非法占有目的，盗伐他人所有的林木。\footnote{根据2023年8月15日最高人民法院《关于审理破坏森林资源刑事案件适用法律若干问题的解释》第3条规定，以非法占有为目的，有下列情形之一的，数量较大的，构成盗伐林木罪：（1）未取得采伐许可证，擅自采伐国家、集体或者他人所有的林木的；（2）违反森林法第五十六条第三款的规定，擅自采伐国家、集体或者他人所有的林木的；（3）在采伐许可证规定的地点以外采伐国家、集体或者他人所有的林木的。不以非法占有为目的，违反森林法的规定，进行开垦、采石、采砂、采土或者其他活动，造成国家、集体或者他人所有的林木毁坏，符合刑法第二百七十五条规定的，以故意毁坏财物罪定罪处罚。}

例1，甲带着非法占有目的，雇佣乙盗伐林木，对乙谎称是自己的林木。乙无罪，甲构成盗伐林木罪的间接正犯。

例2，甲对一片林木有采伐许可证，先口头约定将该林木卖给乙，收取定金，然后又将该林木高价卖给丙，收取了全部货款。在乙向法院起诉甲期间，丙砍伐了该林木。丙拥有所有权，所以不构成盗伐林木罪，也不构成滥伐林木罪。甲对乙构成民事违约。

2.行为对象。一般而言，盗伐活着的、成排成片的林木，才能达到破坏生态环境的程度。

根据司法解释规定，盗伐下列对象，不构成盗伐林木罪，因为这些行为没有破坏生态环境，但是构成盗窃罪：

（1）盗窃他人所有并已经伐倒的树木。

（2）盗伐他人房前屋后、自留地种植的零星树木。（2017年第20题）

（3）盗伐枯死的林木。

3.行为方式。常见方式是砍伐，但以下方式也属于盗伐：

（1）将树木整体挖走，移植到自已地盘。虽然移植后树木没死，但是这种移植行为破坏了原林木所在地的生态环境。（2017年第20题）

（2）剥树皮、掘根等行为，导致树木死亡。如果这些行为没有导致树木死亡，则不构成盗伐林木罪。若有非法占有目的，并且价值数额较大，则构成盗窃罪。如果带着非法占有目的，剥树皮，价值数额较大，并且导致树木死亡，则同时触犯盗窃罪和盗伐林木罪，法条竞合，定盗伐林木罪。（2025年试题）

4.转化抢劫罪。盗伐林木时，为了窝藏赃物、抗拒抓捕或毁灭罪证，使用暴力的，转化为抢劫罪。例如，本节开头“打护林员案”中，甲构成转化型抢劫罪（致人重伤）。

\subsubsection*{（二）滥伐林木罪}

1.行为公式：滥伐林木罪=B（破坏生态环境，没有非法占有目的）。

% Source PDF physical page 384; see audit/source_coverage.json.
2.行为类型：\footnote{根据2023年8月15日最高人民法院《关于审理破坏森林资源刑事案件适用法律若干问题的解释》第5条规定，下列行为属于滥伐林木：（1）未取得采伐许可证，或者违反采伐许可证规定的时间、地点、数量、树种、方式，任意采伐本单位或者本人所有的林木的；（2）违反森林法第五十六条第三款的规定，任意采伐本单位或者本人所有的林木的；（3）在采伐许可证规定的地点，超过规定的数量采伐国家、集体或者他人所有的林木的。林木权属存在争议，一方未取得采伐许可证擅自砍伐的，以滥伐林木论处。}

（1）没有采伐许可证而擅自采伐林木；

（2）虽有采伐许可证，但违反规定采伐林木。

3.行为对象，既包括自己所有的林木，也包括他人所有的林木。滥伐枯死的林木，不构成滥伐林木罪。

4.滥伐林木罪和故意毁坏财物罪可以想象竞合。

司法解释规定：没有非法占有目的，进行开垦、采石、采砂、采土或者其他活动，造成国家、集体或者他人所有的林木毁坏，构成故意毁坏财物罪。

例如，甲为了种植沉香，将大片国有林木砍倒后扔在旁边，然后种植沉香。甲没有采伐证而擅自采伐，构成滥伐林木罪。甲没有非法占有目的，也即没有利用意思，不构成盗窃罪和盗伐林木罪。甲同时构成故意毁坏财物罪。甲的滥伐林木罪和故意毁坏财物罪属于想象竞合，择一重罪论处。（2020年主观题）\footnote{张明楷教授认为，认定甲仅触犯滥伐林木罪和故意毁坏财物罪，滥伐林木罪仅评价了生态环境受到破坏，故意毁坏财物罪仅评价了普通财物受到毁坏，而国有林木所有权受到侵害，被遗漏评价。因此，应认为甲同时也触犯了盗伐林木罪，基于这一考虑，盗伐林木罪不要求具有非法占有目的。由于司法解释要求盗伐林木罪应有非法占有目的，因此张明楷教授的这种观点只能被视为少数说。也即，在主观题开放式的答题中，答哪种观点都给分，但若在客观题中问唯一正确答案，则只能按照司法解释的观点作答。参见张明楷：《盗伐、滥伐林木罪的重要问题》，载《上海政法学院学报》2021年第5期，第10页。}

\subsection*{二、危害珍贵、濒危野生动物罪（第341条）}

1.行为对象。“珍贵、濒危野生动物”，包括某些人工繁殖的动物，例如，人工繁殖的大熊猫。但以食用为目的人工大量繁殖的动物，不是本罪的对象。

2.行为方式。第一，非法猎捕、杀害。第二，非法收购、运输、出售，包括动物制品。

（1）猎捕。例1，甲出于娱乐动机猎捕动物后立即释放的，不构成本罪。例2，甲猎捕动物后取出部分器官，然后放回大自然，不属于杀害，但属于猎捕，构成本罪。

（2）收购。收购动机在所不问，可以为了销售、为了自用。

（3）运输。注意，没有法益侵害的运输行为。例1，动物园管理者未经林业部门批准，将发情老虎送往外地交配的，没有法益侵害，故不构成本罪。例2，居民搬家时运输祖传或年代久远的野生动物制品，没有法益侵害，故不构成本罪。

（4）出售，包括“以营利为目的的加工利用行为”。这是司法解释作的扩大解释。

3.罪数

（1）五个动词都实施，只定一个完整罪名。

（2）使用爆炸、投毒、电网，猎捕杀害，同时危害公共安全的，构成爆炸罪、投放危险物质罪、以危险方法危害公共安全罪，同时构成本罪，想象竞合，择一重罪论处。

（3）先猎捕杀害，然后走私出境，或者先走私人境，然后杀害的，定本罪和走私珍贵

% Source PDF physical page 385; see audit/source_coverage.json.
动物罪，并罚。

（4）直接向走私者收购国家禁止进出口的珍贵动物，不定本罪，而定走私珍贵动物罪。

\subsection*{三、污染环境罪（第338条）}

这是指，违反国家规定，排放、倾倒或者处置有放射性的废物、含传染病病原体的废物、有毒物质或者其他有害物质，严重污染环境的行为。

1.主观要件是故意。

2.既遂标准：严重污染了环境。例如，甲乙各有一套生产设备（电镀槽），共用一个外排口。甲乙都违法排放，严重污染了环境，但不能查明是其中哪套设备严重污染了环境。由于共用一个外排口，二因一果，都构成本罪既遂。

3.实施本罪行为，同时触犯投放危险物质罪的，属于想象竞合犯，择一重罪论处。

4.无危险废物经营许可证从事收集、贮存、利用、处置危险废物经营活动，严重污染环境的，构成污染环境罪，同时构成非法经营罪，想象竞合，择一重罪论处。\footnote{2026年8月15日最高人民法院、最高人民检察院《关于办理环境污染刑事案件适用法律若干问题的解释》。}

\section{毒品犯罪}

\subsection*{一、走私、贩卖、运输、制造毒品罪（第347条）}

\subsubsection*{（一）走私毒品}

走私毒品，是指非法运输毒品进出国境的行为。

1.间接走私。直接向走私毒品的犯罪人购买毒品的，属于走私毒品。

2.既遂标准：陆路输人的，越过国境线；海运、空运的，船舶、飞机到达本国港口、机场。此后被海关查获，不影响既遂的成立。

\subsubsection*{（二）贩卖毒品}

1.贩卖，是指有偿转让。

（1）如果无偿转让毒品，如赠与毒品，不构成贩卖毒品。

（2）有偿转让不要求有牟利目的。例如，甲戒毒后，将自己剩余毒品以极低价格转让给乙，不想赚乙的钱。甲构成贩卖毒品罪。（2008年第65题）

（3）贩卖的方式，可以是钱物交易，也可以是物物交换。例如，甲有海洛因，乙有冰毒。双方交换。甲是吸毒者，想吸食冰毒，想尝尝鲜。乙是贩毒者，想贩卖冰毒，不想吸食海洛因，而是想多个贩卖毒品的品种。乙构成贩卖毒品罪。甲收到10克冰毒后持有，构成非法持有毒品罪。（2024年试题）

（4）为了自己吸食而购买毒品的，不构成贩卖毒品，也不构成贩毒者的贩卖毒品罪、运输毒品罪的共犯。例如，甲为了自已吸食毒品，向乙购买了毒品，给乙发送了收件地址。乙顺利将毒品邮寄到甲家。乙构成贩卖、运输毒品罪。甲不构成共犯。（2022年试题)。

% Source PDF physical page 386; see audit/source_coverage.json.
2.代购算不算贩卖？根据司法解释规定：

\InsertDiagram{figures/manual_p386_1.tex}

[提示]上述图表中要求牟利目的，是为了与友情代购相区分。

例1，甲为了吸食，让乙代购毒品，乙在交通、食宿等必要开销之外，向甲收取“介绍费”“劳务费”。乙从中牟利，变相加价，卖给甲，属于中间商（二道贩子），构成贩卖毒品罪。（2016年第61题）

例2，甲为了吸食，让乙代购毒品，乙说：“我为你代购，不收介绍费、劳务费，但代购来的毒品，分一成给我，作为酬劳。”甲答应。乙其实想将收取的毒品予以贩卖。乙具有贩卖目的，收取毒品，等于是为自己贩卖毒品而进货，构成贩卖毒品罪。乙如果收取毒品是为了自己吸食，则不构成贩卖毒品罪。（2024年试题）

3.贩卖毒品罪的既遂标准是卖掉毒品，不要求收到货款。例如，甲向乙买毒品，给乙交付了首款，乙交付了毒品。甲在ATM机上转尾款时，甲乙被抓获。乙构成贩卖毒品罪既遂。（2024年试题）

4.共犯问题总结：

\InsertDiagram{figures/manual_table_p386_2.tex}

% Source PDF physical page 387; see audit/source_coverage.json.
\ContinuedTable{figures/manual_table_p387_1.tex}

\subsubsection*{（三）运输毒品}

运输毒品，是指使毒品发生较长距离的位移。如果是从一楼搬到三楼，不构成运输毒品罪。

1.运输毒品罪与非法持有毒品罪的关系

[总结]司法解释规定：\footnote{2023年6月29日最高人民法院《全国法院毒品案件审判工作会议纪要》。}（2022年试题）

\InsertDiagram{figures/manual_table_p387_2.tex}

例1，甲为了吸食，买到10克海洛因放在家里，属于持有数量较大，构成非法持有毒品罪。

注意：为了吸食的购买行为本身不构成犯罪。

例2，乙为了吸食，携带10克海洛因出差坐火车，属于运输数量较大。乙构成运输毒品罪。

提示：原本，只有为了走私、贩卖而运输毒品的，才能定运输毒品罪；为了吸食而运输毒品的，应定非法持有毒品罪。但是，由于主观目的不容易查证，所以司法解释就“一刀切”，在认定运输毒品罪时，不考虑主观目的是为了走私、贩卖，还是为了吸食，只考虑是否在实施运输行为。

例3，丙为了吸食，携带5克海洛因出差坐火车，由于运输数量不大，不构成运输毒品罪，也不构成非法持有毒品罪。

［注意1］数量问题。普通的运输毒品罪的成立，不要求运输数量较大。而为了吸食，运输毒品，要成立运输毒品罪，要求运输数量较大。

[注意2］共犯问题。例如，甲乙在某地出差，看到海洛因较便宜，商量买点回去吸食。甲买了6克，乙买了4克，各自购买，各自保管，一起坐大巴车回家。（1）甲乙在购买环节有共同商量，但是购买行为本身不构成犯罪，所以“共同商量购买”这个环节不构成共同犯罪。（2）在保管及运输环节，由于是各自保管，所以也不构成共同犯罪，只能单独处理。由于各自运输数量不大，均未达到10克，甲乙不构成运输毒品罪和非法持有毒品罪。（2025年试题）如果甲交给乙一起保管，乙保管，二人一起坐大巴车回家，那么，乙的运输数量达到10克，构成运输毒品罪的实行犯，甲构成该罪的共犯。

2.运输毒品罪的着手与既遂

（1）着手标准：为了运输而开始搬运。

% Source PDF physical page 388; see audit/source_coverage.json.
（2）既遂标准：毒品离开原存放地，发生较长距离的位移。

[注意1]既遂的标准不是运抵目的地。中途被查获，仍构成既遂。（2017年第61题）

[注意2］行为人先将毒品从甲地运至乙地，由于某种原因，又运回甲地的，构成运输毒品罪既遂。（2022年试题）

3.共犯问题。

第一，分段运输。老大安排，由甲走第一段路程，然后将货交给乙，乙走第二段路程。虽然甲乙有交接，但是二人不构成运输毒品罪的共同犯罪，因为大家都是马仔，各自负责各自的事情。（2016年第61题）

第二，同行运输。老大安排，甲乙各自携带毒品，同行运输毒品，路上甲乙的运输行为相对独立，没有配合、掩护行为，各自向老大领取报酬，不构成运输毒品罪的共同犯罪。

第三，集中派活。老大将甲乙叫到跟前，吩咐甲送一批货到A目的地，乙送一批货到B目的地。甲乙虽然知道对方都在运输毒品，但不构成共同犯罪，因为大家都是马仔，各自负责各自的事情。

\subsubsection*{(四）制造毒品}

1.制造方法

（1）从原料中提取。例如，从罂粟中提取鸦片。

（2）精制，也即提高纯度。例如，对海洛因进行浓缩精华，提高纯度。注意，去除毒品中的非毒品物质，不属于制造毒品。

（3）变种，也即用化学方法或其他方法使一种毒品变成另一种毒品。

（4）非法按照一定的处方针对特定人调制毒品。

[注意1］为便于隐蔽运输、销售、使用、欺骗购买者，或者为了增重，对毒品掺杂使假，添加其他非毒品物质，不属于制造毒品。（2010年第60题）

[注意2]分装毒品属于制造毒品。这是指将毒品进行分割，并装入一定的容器，属于量的精制。

2.既遂标准：制造出来毒品。至于制造出来的毒品数量多少、纯度高低，都不影响既遂的成立。（2022年试题）

［注意］不能犯。制造毒品的方法在客观上根本不具有制造出毒品的可能性的，不成立制造毒品罪，注意不能以制造毒品罪未遂论处。

\subsubsection*{（五）主观认识问题}

只要求认识到毒品的成分作用，不要求认识到毒品名称。

例1，张某发现某种电子烟中的大麻素能够使人精神成瘾，但不知道其是“毒品”，仍出售这种电子烟。毒品是指能够使人精神成瘾的精神药品。张某已经认识到这种成分作用，便认识到毒品的实质特征，具有贩卖毒品故意。张某构成贩卖毒品罪。（2022年试题）

例2，甲误以为自己贩卖的是海洛因，实际是鸦片。要具有贩卖毒品罪的故意，只要求认识到“毒品”即可，不要求认识到具体的毒品的种类。甲构成贩卖毒品罪。（2008年第2题）

% Source PDF physical page 389; see audit/source_coverage.json.
\subsubsection*{（六）麻精药品（麻醉药品、精神药品）\protect\footnotemark}\footnotetext{2023年6月26日最高人民法院《全国法院毒品案件审判工作会议纪要》。}

1.明知是走私、贩卖毒品的犯罪分子或者吸毒者，向其贩卖麻精药品，构成贩卖毒品罪。

2.出于医疗目的，未取得药品证件，生产、进口、销售麻精药品，不构成毒品犯罪，构成妨害药品管理罪。

3.出于为自己治病目的，携带、邮寄麻精药品进出境，不构成毒品犯罪。

4.依法管理麻精药品的人员：（1）出于医疗目的，未经许可，非法贩卖麻精药品，不构成毒品犯罪，多数说认为定非法经营罪。\footnote{参见喻海松编著：《实务刑法评注》（第三版），北京大学出版社2025年版，第2026页。}（2）向吸毒者无偿提供麻精药品，构成非法提供麻醉药品、精神药品罪（第355条）。（3）向吸毒者有偿提供（贩卖）麻精药品，构成非法提供麻醉药品、精神药品罪和贩卖毒品罪，想象竞合，择一重罪论处，定贩卖毒品罪。（4）明知是走私、贩卖毒品的犯罪分子，向其贩卖麻精药品，构成贩卖毒品罪；向其无偿提供麻精药品，属于供货行为，构成走私、贩卖毒品罪的帮助犯。

\subsubsection*{（七）数量问题}

成立本罪，不要求毒品数量大小。但数量影响量刑。\CJKunderwave{数量计算问题。}

1.走私、贩卖、运输、制造的毒品，如果是同一宗毒品，在数量上不重复计算；如果是不同宗毒品，在数量上需要累加计算。不过无论是否同宗毒品，在定罪上都只定走私、贩卖、运输、制造毒品罪，不数罪并罚，因为该罪名是选择性罪名。（2010年第60题）

2.走私、贩卖、运输、制造、非法持有两种以上毒品的，可以将不同种类的毒品分别折算为海洛因的数量，以折算后累加的毒品总量作为量刑的根据。（2019年试题）

\subsubsection*{（八）特别再犯制度}

第356条因走私、贩卖、运输、制造、非法持有毒品罪被判过刑，又犯本节规定之罪的，从重处罚。

1.前罪只包括五个罪名：走私、贩卖、运输、制造、非法持有毒品罪。后罪包括毒品犯罪一节所有的罪名。

2.前后罪的法定刑没有要求，包括管制、拘役、有期徒刑、无期徒刑。

3.前后罪的时间间隔没有要求。

4.特别再犯与累犯的关系。

（1）二者可以竞合。行为人既构成累犯，又构成特别再犯，产生竞合，择一重罚处罚。

(2）不满18周岁的人不构成累犯，也不构成特别再犯。（2013年第3题)

\subsection*{二、非法持有毒品罪（第348条）}

1.成立本罪要求持有毒品达到数量较大。例如，海洛因是10克（2011年第18题）。甲乙各自买了5克海洛因，乙将5克海洛因交给甲保管。甲共保管了10克海洛因。甲构成本罪，乙构成本罪的共犯。

2.持有是指事实上支配控制毒品。

（1）持有不限于随身携带，也包括保存在可以控制的地方。

（2）持有也可以是间接持有。例如，委托他人代为保管。委托者属于间接持有，成立

% Source PDF physical page 390; see audit/source_coverage.json.
非法持有毒品罪。代为保管人明知是毒品而代为保管，也成立非法持有毒品罪。如果代为保管人同时还具有帮助原持有人逃避司法机关追查的意图，则既触犯非法持有毒品罪，还触犯窝藏毒品罪，属于想象竞合，择一重罪论处。（2011年第18题，2010年第60题）

3.因走私、贩卖、运输而持有毒品，定走私、贩卖、运输毒品罪，不定非法持有毒品罪。（2011年第18题）

4.吸毒不构成犯罪，但是吸毒者持有毒品达到数量较大时，以非法持有毒品罪论处。

\subsection*{三、容留他人吸毒罪（第354条）}

1.容留，是指允许他人在自己管理的场所吸食、注射毒品或者为他人吸食、注射毒品提供场所的行为。容留行为既可以主动实施，也可以被动实施；既可以是有偿的，也可以是无偿的。

2.共有场所。如果二人共有某个场所，其中一人在吸毒，另一人不阻止，不属于容留他人吸毒。

3.罪数问题。根据司法解释，\footnote{2016年4月11日最高人民法院《关于审理毒品犯罪案件适用法律若干问题的解释》。}向他人贩卖毒品后又容留其吸食、注射毒品，或者容留他人吸食、注射毒品并向其贩卖毒品，以贩卖毒品罪和容留他人吸毒罪数罪并罚。（2017年第61题）

\subsection*{四、引诱、教唆、欺骗他人吸毒罪和强迫他人吸毒罪（第353条）}

例如，店家将罂粟壳掺入食品以招揽顾客。由于顾客都是在不知道的情况下被骗食用的，因此这种行为属于欺骗他人吸食毒品的犯罪行为。

\subsection*{五、包庇毒品犯罪分子罪和窝藏、转移、隐瞒毒品、毒赃罪（第349条）}

\subsubsection*{（一）包庇毒品犯罪分子罪}

1.包庇的对象不是任何毒品犯罪分子，而只是走私、贩卖、运输、制造毒品罪的犯罪分子。（2012年第62题，2021年试题）

2.本罪与包庇罪是法条竞合关系，优先适用本罪。

\subsubsection*{（二）窝藏、转移、隐瞒毒品、毒赃罪}

本罪属于事后的赃物犯罪。本罪是掩饰、隐瞒犯罪所得罪的一个具体罪名，二者是法条竞合关系，优先适用本罪。（2012年第62题，2021年试题）

1.本罪的上游犯罪，不包括所有的毒品犯罪，只包括“走私、贩卖、运输、制造毒品罪”。

2.本罪只能发生在上游犯罪既遂之后。若与上游犯罪人事前有通谋，则构成上游犯罪的共同犯罪。

例1，甲走私毒品既遂后，为逃避法律追究，让乙为其窝藏该批毒品。乙照办。乙构成窝藏毒品罪。如果甲与乙事前通谋，由甲走私入关，由乙负责窝藏，则乙构成走私毒品罪的共犯。

例2，甲贩卖毒品既遂后，将获得的货款（也即毒赃）交给乙窝藏，乙照办。乙构成

% Source PDF physical page 391; see audit/source_coverage.json.
窝藏毒赃罪。如果甲与乙事前通谋，由甲贩卖，由乙窝藏获得的货款，则乙构成贩卖毒品罪的共犯。

例3，甲将毒品从上海运到南京后，为了逃避警方侦查、法律追究，将毒品交给乙，让乙窝藏或转移到秘密地点，乙照办。乙构成窝藏、转移毒品罪。如果甲与乙事前通谋，由甲将毒品从上海运到南京，由乙将毒品从南京再运到北京，则甲、乙构成运输毒品罪的共同犯罪。

3.由于本罪是事后的赃物犯罪，因此，实施本罪的目的是逃避法律追究。如果直接目的就是贩卖，则直接定贩卖毒品罪。

例1，甲走私毒品、运输毒品、制造毒品既遂后，为了贩卖，将毒品交给乙，让乙贩卖掉。乙答应，乙先窝藏、转移，然后找机会贩卖掉。乙直接构成运输、贩卖毒品罪，不认定为窝藏、转移毒品罪，也不认定为掩饰、隐瞒犯罪所得罪。

例2，甲将毒品交给乙，让乙运到北京，并卖给吸毒者。乙照办。乙构成运输、贩卖毒品罪，不构成转移毒品罪。（2015年第9题）

4.行为对象：“毒品和毒赃”。“毒品”，是指走私、运输、制造毒品罪既遂后的毒品。贩卖毒品罪既遂，就意味着卖掉了毒品，换来货款，所以不存在“毒品”，只存在“毒赃”（获得的货款，赃款）。

[注意]本罪的“毒赃”，不包括“毒赃”产生的收益。例如，甲通过贩毒获得50万元货款（毒赃），用这笔钱炒股赚了10万元（毒赃的收益）。乙帮甲保管这10万元。乙不构成窝藏毒赃罪，但构成掩饰、隐瞒犯罪所得收益罪。

\subsection*{六、非法种植毒品原植物罪（第351条）}

本罪是指，非法种植罂粟、大麻等毒品原植物，情节严重的行为。具有下列三种情形之一的，属于“情节严重”：（1）种植数量较大。（2）经公安机关处理后又种植。（3）抗拒铲除。

本罪的“种植”，包括播种、育苗等。本罪的既遂标准：只要种下数量较大的种子，就构成既遂（司法解释）。不过，有时候，数量不大，但满足后两种“情节严重”的，也值得处罚。

例如，甲种植大量罂粟种子，种子刚发芽即被大水冲垮。甲构成本罪既遂。（2025年试题）

\OriginalPanel{典型真题}

\ensuremath{>}\ensuremath{>}

关于毒品犯罪，下列哪些选项是正确的？（2017年·卷二·61题)\footnote{[答案]ABCD。}

A.甲容留未成年人吸食、注射毒品，构成容留他人吸毒罪

B.乙随身携带藏有毒品的行李入关，被现场查获，构成走私毒品罪既遂

C.丙乘广州至北京的火车运输毒品，快到武汉时被查获，构成运输毒品罪既遂

D.丁以牟利为目的容留刘某吸食毒品并向其出卖毒品，构成容留他人吸毒罪和贩卖毒品罪，应数罪并罚

% Source PDF physical page 392; see audit/source_coverage.json.
\section{卖淫类犯罪}

[问题］卖淫女乙嫌老板甲给的待遇太低，想跳槽到另一场所从事卖淫。甲不同意，强迫乙继续在自已的场所卖淫。甲是否构成强迫卖淫罪？（禁止“人才”流动案）

\subsection*{一、组织卖淫罪、强迫卖淫罪、协助组织卖淫罪（第358条）}

\subsubsection*{（一）卖淫}

1.卖淫者的范围。由于组织、强迫、引诱、容留、介绍卖淫罪中，规定的卖淫者是“他人”，所以既包括由女性向男性卖淫，也包括由男性向女性卖淫；既包括由女性向女性卖淫，也包括由男性向男性卖淫。（2009年第54题）

2.接受卖淫服务的人必须是不特定人。组织女性被某个特定人包养的，不应认定为组织卖淫罪。

3.卖淫的支付方式。卖淫是钱色交易，钱指金钱或有价值财物。权色交易不算卖淫。同理，帮助女方办某件事（如写作业），女方提供性服务，不算卖淫。

\subsubsection*{（二）组织卖淫罪}

1.组织卖淫，要求对卖淫者具有管理性和控制性。这是组织卖淫与介绍卖淫的区别。有无固定场所，在所不问。

2.被组织者实施了卖淫行为的，组织卖淫罪构成既遂。

\subsubsection*{（三）强迫卖淫罪}

1.逼良为。在他人不愿意从事卖淫活动的情况下，使用强制手段迫使其从事卖淫活动。

2.不让从良。他人虽然原本从事卖淫活动，但在他人不愿意继续从事卖淫活动的情况下，使用强制手段迫使其继续从事卖淫活动。

3.具体时间、地点的强迫。例如，本节开头“禁止‘人才’流动案”中，甲在乙不愿意在此地从事卖淫活动的情况下，强迫乙继续卖淫，构成强迫卖淫罪。

[注意1]强迫他人与特定的个人性交或者从事猥亵活动的，成立强奸、强制猥亵等罪。

[注意2】行为人强迫妇女仅与自己发生性交，并支付性行为对价的，应认定为强奸罪，不得认定为强迫卖淫罪。

4.被强迫者实施了卖淫行为的，强迫卖淫罪构成既遂。

\subsubsection*{（四）协助组织卖淫罪}

1.行为方式：招募、运送卖淫者，充当保镖、打手、管账人等。

[注意］在具有营业执照的会所、洗浴中心等经营场所担任保洁员、收银员、保安员等，从事一般服务性、劳务性工作，仅领取正常薪酬，不认定为协助组织卖淫罪。\footnote{2017年7月25日最高人民法院、最高人民检察院《关于办理组织、强迫、引诱、容留、介绍卖淫刑事案件适用法律若干问题的解释》。}

2.本罪是将帮助行为作为实行行为对待，将帮助行为正犯化，成为独立罪名。例如，乙告知甲打算组织卖淫，让甲物色一些“人才”。甲便在市场上公开招募卖淫女，10名妇女知道从事卖淫，应召前来。甲送10名妇女给乙的路上被抓。虽然乙尚未实施组织卖淫，

% Source PDF physical page 393; see audit/source_coverage.json.
但甲的行为具有法益侵害性，值得刑罚处罚，构成协助组织卖淫罪。

\subsection*{二、引诱、容留、介绍卖淫罪和引诱幼女卖淫罪（第359条）}

1.行为方式。

（1）引诱，是指在他人本无卖淫意愿的情况下，使用勾引、利诱等手段使他人从事卖淫活动的行为。卖淫者原本在此时、此地卖淫，行为人引诱其在彼时、彼地卖淫的，不应认定为引诱他人卖淫。

（2）容留，是指允许他人在自己管理的场所卖淫或者为他人卖淫提供场所的行为。

（3）介绍，是指在卖淫者与客之间牵线搭桥，沟通撮合，使他人卖淫得以实现的行为。

[注意1]在意欲卖淫者与卖淫场所的管理者之间进行介绍的，属于介绍他人卖淫。但是，单纯向意欲者介绍卖淫场所，而与卖淫者没有任何联络的，可谓“介绍他人”，不能认定为介绍卖淫。介绍女子被他人“包养”的，不成立介绍卖淫罪。

[注意2］引诱、容留、介绍他人卖淫是否以营利为目的，不影响犯罪的成立。

[注意3]引诱与强奸。例如，甲引诱4名13周岁的幼女卖淫，并让她们对外称16周岁。客乙不知情，和幼女发生性关系。乙不构成强奸罪。甲构成引诱幼女卖淫罪。甲同时构成强奸罪的间接正犯，因为甲欺骗乙，利用乙的不知情，间接侵害了幼女的法益，构成奸淫幼女型强奸罪。甲的两罪想象竞合，择一重罪论处。（2024年试题）

2.总结罪数关系：（2004年第89题）

\InsertDiagram{figures/manual_table_p393_1.tex}

3.总结常考的引诱型罪名：

（1）引诱、容留、介绍卖淫罪（第359条第1款）。

（2）引诱幼女卖淫罪（第359条第2款）。

（3）引诱未成年人聚众淫乱罪（第301条第2款。\CJKunderwave{注意}：比照聚众淫乱罪从重处罚）。

（4）引诱他人吸毒罪（第353条第1款。\CJKunderwave{注意}：引诱未成年人吸毒不是独立罪名，只是引诱他人吸毒罪的从重处罚情节）。

\subsection*{三、传播性病罪}

第360条明知自己患有梅毒、淋病等严重性病\CJKunderwave{卖淫、嫖娼的}，处五年以下有期徒刑、拘役或者管制，并处罚金。

1.行为方式：严重性病患者卖淫、。本罪是行为犯（抽象危险犯），只要严重性

% Source PDF physical page 394; see audit/source_coverage.json.
病患者卖淫、，就构成本罪。至于是否实际造成他人患上严重性病的后果，不影响本罪的成立。

2.主观上要求明知自己患有严重性病。具有下列情形的，应认定为“明知”：第一有证据证明曾到医院或者其他医疗机构就医或者检查，被诊断为患有严重性病的；第二根据本人的知识和经验，能够知道自己患有严重性病的。

3.本罪中的性病包括艾滋病。这是司法解释作的扩大解释。\footnote{2017年7月25日最高人民法院、最高人民检察院《关于办理组织、强迫、引诱、容留、介绍卖淫刑事案件适用法律若干问题的解释》。}

4.意图伤害他人，通过卖淫、使对方染上严重性病，既构成本罪，又构成故意伤害罪，择一重罪论处，定故意伤害罪。其中，使对方感染上艾滋病病毒的，构成故意伤害罪（重伤）。例如，明知自己感染艾滋病病毒而卖淫、的；明知自己感染艾滋病病毒，故意不采取防范措施而与他人发生性关系，使对方感染艾滋病的，以故意伤害罪（重伤）论处。

\section{关于淫秽物品的犯罪}

[问题］甲女在网络直播间进行淫秽表演，吸引观众打赏。甲是否构成传播淫秽物品牟利罪？（淫秽直播案）

\subsection*{一、制作、复制、出版、贩卖、传播淫秽物品牟利罪（第363条）}

1.“淫秽物品”。根据司法解释\footnote{2004年9月6日最高人民法院、最高人民检察院《关于办理利用互联网、移动通讯终端、声讯台制作、复制、出版、贩卖、传播淫秽电子信息刑事案件具体应用法律若干问题的解释（一）》。}规定，存储在电脑硬盘、网络云盘、网络服务器中的淫秽视频文件（电子信息）属于淫秽物品。网上裸聊、网上淫秽直播属于即时存储在网络服务器中的电子信息，属于淫秽物品。例如，本节开头的“淫秽直播案”中，甲构成传播淫秽物品牟利罪。

2.贩卖：是指有偿转让淫秽物品。\CJKunderwave{注意}：以出售为目的的购买行为是贩卖的预备行为，以传播为目的的购买行为是传播的预备行为。以供自己使用的购买行为不构成犯罪。（2012年第55题）

3.传播：是指让不特定人或多数人感知。传播的途径有两种：一是主动传播出去，二是\CJKunderwave{陈列}，等候不特定人主动来看。也即，陈列也是传播的一种方式。

例1，某夜总会在临街橱窗陈列淫秽书画，招揽顾客，属于陈列淫秽物品，属于传播淫秽物品。

例2，某夜总会让几名裸体卖淫女站在临街橱窗，招揽顾客，虽是陈列，但不是陈列淫秽“物品”，不属于传播淫秽物品。

[提示］传播与贩卖不是对立排斥关系，有些行为既是传播，也是贩卖。例如，在网上提供有偿下载淫秽视频文件的服务，既是传播，也是贩卖。

% Source PDF physical page 395; see audit/source_coverage.json.
\subsection*{二、传播淫秽物品罪（第364条）}

[总结】三个罪名中的“目的”

\InsertDiagram{figures/manual_table_p395_1.tex}

\subsection*{三、大总结：违禁品的罪数}

1.盗窃+销售

\InsertDiagram{figures/manual_table_p395_2.tex}

[注意］对象认识错误。例如，甲欲盗窃普通财物，到手后发现是枪支，又出售。盗窃行为定盗窃罪，出售行为定非法买卖枪支罪，数罪并罚。

2.走私+销售

\InsertDiagram{figures/manual_table_p395_3.tex}

\subsection*{四、大总结：刑法中的“贩卖”“销售”“出售”“倒卖”“买卖”}

1.共同点：客观上要求有偿转让。

第一，如果是无偿转让，则不是卖，而是送。

第二，有偿的方式，不限于金钱货币，可以是其他物质利益，可以物物交换。

例1，甲给乙1克毒品，乙给甲10袋大米，也即物物交换也是买卖。

% Source PDF physical page 396; see audit/source_coverage.json.
例2（2014年第57题），甲用毒害性物质，乙用放射性物质，二者交换。甲、乙的行为属于非法买卖危险物质罪。

例3，甲用枪，乙用子弹，二者交换。二者构成非法买卖枪支、弹药罪。

第三，主观上不要求有营利目的。也即，赔本卖，也是卖。

2.出卖行为：“贩卖”“销售”“倒卖”。

第一，这里的出卖，包括先买进后卖出，也包括单纯的卖出。注意，倒卖文物罪中的“倒卖”以前要求先买进后卖出，新的司法解释规定，没有买进行为，只有卖出行为也构成倒卖文物罪。

第二，这些罪名不处罚为了自用的购买行为。对于为了出售而购买的行为，一般是作为出售的预备行为来处理。例如，为了贩卖而购买毒品，作为贩卖毒品罪的预备行为。

3.“买卖”。

既处罚“卖”，也处罚“买”。卖，包括先买进后卖出，也包括单纯的卖出。买，包括为卖出而买进，也包括为自已使用而买进。例如，非法买卖枪支、弹药罪，非法买卖危险物质罪，买卖国家机关公文、证件、印章罪，买卖身份证件罪，这些罪中的“买卖”都应如此理解。

例1（2014年第57题），吸毒者甲用毒害性物质，贩毒者乙用毒品，二者交换。甲出卖毒害性物质，构成非法买卖危险物质罪。乙购买毒害性物质，构成非法买卖危险物质罪。甲购买毒品自用，不构成贩卖毒品罪。乙出卖毒品，构成贩卖毒品罪。乙貌似实施了一个行为，实际实施了两个行为，购买行为和出卖行为，因此应数罪并罚。

例2，吸毒者甲用国家机关公文，贩毒者乙用毒品，二者交换。甲构成买卖国家机关公文罪。乙构成贩卖毒品罪和买卖国家机关公文罪，并罚。

4.既遂条件。

第一，出售是既遂条件。例如，生产、销售伪劣产品罪。只有销售金额达到5万元，犯罪才既遂。如果仅有生产行为，没有销售出去，库存额有15万元，构成犯罪未遂。

第二，出售不是既遂的条件。

例1（2018年试题），制作、复制、出版、贩卖、传播淫秽物品牟利罪。每一个行为都是能够达到既遂的行为，带着牟利目的完成其中一个行为，就构成既遂。构成既遂，不要求将制作、复制的淫秽物品贩卖出去。

例2，走私、贩卖、运输、制造毒品罪。每一个行为都是能够达到既遂的行为，完成其中一个行为，就构成既遂。构成既遂，不要求将制造的毒品贩卖出去。

例3，出售、购买、运输假币罪。每一个行为都是能够达到既遂的行为，完成其中一个行为，就构成既遂。构成既遂，不要求将购买的假币出售出去。
